\documentclass[11pt,twocolumn]{article}

% ─── Packages ──────────────────────────────────────────────
\usepackage[utf8]{inputenc}
\usepackage[T1]{fontenc}
\usepackage{amsmath,amssymb,amsthm}
\usepackage{mathtools}
\usepackage{bm}
\usepackage{graphicx}
\graphicspath{{figures/}{./}}
\usepackage{booktabs}
% \usepackage{algorithm}     % unused; re-enable if you add algorithm envs
% \usepackage{algorithmic}   % unused; re-enable if you add algorithmic envs
\usepackage[margin=1in]{geometry}
\usepackage{hyperref}
\usepackage{cleveref}
\usepackage{natbib}
\usepackage{xcolor}
% \usepackage{thmtools}   % unused; re-enable if you add \declaretheorem
\usepackage{caption}

% ─── Theorem environments ──────────────────────────────────
\newtheorem{theorem}{Theorem}[section]
\newtheorem{lemma}[theorem]{Lemma}
\newtheorem{proposition}[theorem]{Proposition}
\newtheorem{corollary}[theorem]{Corollary}
\newtheorem{definition}[theorem]{Definition}
\newtheorem{remark}[theorem]{Remark}

% ─── Custom commands ───────────────────────────────────────
\newcommand{\R}{\mathbb{R}}
\newcommand{\norm}[1]{\left\lVert #1 \right\rVert}
\newcommand{\abs}[1]{\left| #1 \right|}
\newcommand{\eisv}{\ensuremath{\mathbf{x}}}
\newcommand{\deta}{\ensuremath{\Delta\bm{\eta}}}
\newcommand{\cv}{\ensuremath{C(V,\Theta)}}
\newcommand{\todo}[1]{\textcolor{red}{\textbf{[TODO: #1]}}}

% ─── Metadata ──────────────────────────────────────────────
\title{UNITARES: Information-Theoretic Governance\\
of Heterogeneous Agent Fleets}

\author{
  Kenny Wang\\
  Independent Researcher\\
  \texttt{founder@cirwel.org}
}

\date{Draft --- \today}

% ═══════════════════════════════════════════════════════════
\begin{document}
\maketitle

% ─── Abstract ──────────────────────────────────────────────
\begin{abstract}
UNITARES is a runtime governance layer for heterogeneous AI-agent
fleets. It tracks continuous agent state, calibrates by class,
detects drift, and issues governance interventions with auditable
provenance. Each agent carries a four-dimensional state vector
$\eisv = (E, I, S, V)^\top$ updated on every check-in, intended to
support continuous governance rather than discrete after-the-fact
filtering. Fleet-wide normalization is the wrong target: the
populations we govern---embodied agents with sensor-driven state,
persistent autonomous services, session-bounded coding assistants,
and ephemeral parser agents---do not share an output modality, a
tempo, or a healthy operating point.

UNITARES addresses this with two structural moves. First, it
reinterprets the EISV coordinates in information-theoretic terms:
$S$ as response-distribution entropy, $I$ as context-response
mutual information, $E$ as negative variational free energy, and
$V$ as an accumulated free-energy residual. Second, it makes
normalization class-conditional: scale constants and healthy
operating points are functions of an agent class keyed on
existing identity tags, with fleet-wide defaults retained only as
fallback behavior.

Because UNITARES was already deployed, re-grounding the
coordinates raised a live-systems problem in addition to a
mathematical one. We introduce a pipeline-ordering migration
mechanism that separates decision-producing from
presentation-producing pipeline stages so a deployed governance
framework can acquire new coordinate semantics without breaking
live behavior; we argue this mechanism generalizes to any
deployed agent framework facing a semantic-coordinate change.
Stability of the dynamics is preserved under the reformulation
and established by contraction analysis (\Cref{app:contraction-proof}).

We illustrate the framework with a production deployment snapshot from
February 20, 2026, covering 903 registered agents and 75 active agents.
A verdict counterfactual on a 30-day production slice
($N = 13{,}310$) shows that $28.9\%$ of basin assignments would flip
under class-conditional grounded coherence, with flip rates ranging
$15$--$33\%$ per class and a strong directional bias into the
\texttt{low} basin, showing that the homogenization-failure-mode
argument has first-order consequences at the gating layer rather
than only at the reported-value layer. A first per-class calibration
of the manifold-coherence path is reported here; the remaining scale
constants and higher-tier estimators (logprob and multi-sample) are
deferred to subsequent work.
\end{abstract}

\textbf{Keywords:} autonomous governance, information theory,
variational free energy, heterogeneous multi-agent systems,
class-conditional calibration, AI safety

\vspace{1ex}
\noindent\fbox{\parbox{0.97\columnwidth}{\small
\textbf{Reading aids.} We introduce four terms up front:
\emph{Variational free energy} ($-F$, used as the $E$ coordinate):
the Free-Energy-Principle quantity measuring how well an agent's
predictions match outcomes; high $E$ = low surprise.
\emph{Cosmological soup}: the failure mode of averaging a
heterogeneous agent fleet into a single state distribution, in
which no population structure survives normalization.
\emph{Manifold coherence} ($C = 1 - \norm{\Delta}_2 / \norm{\Delta}_{\max}$):
distance from the agent class's healthy operating point, used as
a class-conditional replacement for the legacy
$\tanh$-of-$V$ coherence.
\emph{Basin flip}: a change in governance-basin assignment
(\texttt{high} / \texttt{boundary} / \texttt{low}) when the same
state vector is classified under different coherence formulas.
}}

% ═══════════════════════════════════════════════════════════
\section{Introduction}
\label{sec:intro}

Autonomous AI agents operating in production environments face a
fundamental tension: they must explore to learn, yet remain safe
while doing so. Existing approaches to agent governance fall into
two camps. \emph{Hard constraints} (guardrails, output filters)
provide safety but impede capability. \emph{Soft monitoring}
(logging, dashboards) preserves capability but offers no formal
guarantees.

We argue that governance should be \emph{intrinsic}: a
continuous dynamical process that shapes agent behavior from
within, producing leading indicators of degradation rather than
binary verdicts after the fact. This paper presents UNITARES, a
framework
in which each agent carries a four-dimensional state
$\eisv = (E, I, S, V)^\top$ whose evolution is governed by coupled
differential equations. The state provides a continuous,
interpretable signal for governance decisions---when to proceed,
when to pause, when to escalate.

UNITARES is a deployed framework; the present formulation is
organized around two structural commitments: information-theoretic
grounding of the EISV coordinates, and class-conditional
calibration that treats heterogeneity as a first-class design
constraint. The populations
the framework governs are heterogeneous---they differ in tempo by
four orders of magnitude, in output modality (text, sensor,
system action), and in healthy operating point. A governance
framework that collapses this population to a single state
distribution---what we call \emph{cosmological soup}---loses
signal exactly where signal matters. The contributions below are
organized around these two commitments.

\subsection{Contributions}

This paper makes five contributions:

\begin{enumerate}
\item \textbf{Framework reframing for heterogeneous fleets}
  (\Cref{sec:heterogeneity}): We identify heterogeneity, rather than
  single-population stability, as the governing design constraint for a
  production agent-governance framework.

\item \textbf{Information-theoretic grounding of EISV}
  (\Cref{sec:grounding}): We give each coordinate computable
  information-theoretic semantics, together with tiered estimation
  recipes and provenance tags.

\item \textbf{Class-conditional calibration protocol}
  (\Cref{sec:class-calibration}): We formalize scale constants and
  healthy operating points as class-indexed quantities derived from
  identity structure, with explicit fallback behavior for sparse or
  unclassified classes.

\item \textbf{Behavior-preserving migration mechanism for deployed systems}
  (\Cref{sec:migration}): We introduce a pipeline-ordering strategy
  that separates decision-producing and presentation-producing stages so
  semantic re-grounding can occur without changing live governance
  behavior during migration.

\item \textbf{Production deployment snapshot and verdict counterfactual}
  (\Cref{sec:experiments}): We report a production snapshot of the
  deployed system to characterize operating regimes, saturation
  behavior, and observed fleet dynamics. We also report a verdict
  counterfactual on a $13{,}310$-row production slice showing that
  $28.9\%$ of basin assignments would flip under class-conditional
  grounded coherence, with the dominant direction of flips
  (predominantly into the \texttt{low} basin) showing that the
  homogenization-failure-mode argument of contribution~$1$ has
  first-order consequences at the gating layer, not only at the
  reported-value layer.
\end{enumerate}

\subsection{Limitations and Scope}
\label{sec:limitations-and-scope}

This paper is a framework-and-migration paper, not a completed
empirical validation paper. Its primary claims are architectural and
methodological: that heterogeneous fleets require class-conditional
calibration, that the EISV coordinates can be reinterpreted
information-theoretically without changing the governing dynamics, and
that a deployed governance system can be re-grounded by separating
decision-producing from presentation-producing stages in the request
pipeline.

Accordingly, three boundaries matter. First, the class-conditional
calibration protocol is specified here, and a first measurement
populating the manifold-coherence path's per-class constants is
reported in \Cref{sec:phase-2-measurement}; the remaining
$S_{\text{scale}}, I_{\text{scale}}, E_{\text{scale}}$ constants
(used only by the higher-fidelity tier-1 and tier-2 estimators)
remain placeholders awaiting those estimators landing. Second, the
highest-fidelity estimators for grounded quantities (logprob-based
and multi-sample computation) remain future instrumentation work;
the current deployment relies on degraded-mode fallbacks where
those signals are unavailable. Third, the deployment results in
\Cref{sec:experiments} are presented as a single descriptive
production snapshot rather than a controlled validation study
across stable parameter regimes.

The paper should therefore be read as establishing the framework, the
migration mechanism, and the shape of the empirical program, while
deferring full class-conditional validation to subsequent work.

\subsection{Related Work}

\paragraph{Information-Theoretic Monitoring.}
Shannon entropy and mutual information~\cite{shannon1948} are standard tools for
characterizing model output distributions; calibration analysis
\cite{guo2017calibration} treats the related question of
confidence-outcome alignment. The variational free-energy
principle~\cite{friston2010free} provides a unifying frame in
which agent surprise corresponds to free energy and adaptive
agents minimize expected surprise; the active-inference literature
\cite{dacosta2020active,sajid2021active,smith2022step,parr2022active}
gives a modern treatment of that lineage. We take from this
literature the semantic referent for the $E$ coordinate (the
variational free-energy concept of~\cite{friston2010free}); we do
not apply the discrete-state-space machinery of the active-inference
treatments, because the agents under governance do not generally
expose discrete latent state-spaces to the governance layer.
UNITARES adapts the remaining information-theoretic quantities
(Shannon entropy, mutual information) to the governance setting
and treats the choice of computation tier (logprob, multi-sample,
heuristic) as a deployment configuration rather than an
assumption.

\paragraph{Contraction Theory.}
Lohmiller and Slotine~\cite{lohmiller1998contraction} established
contraction analysis as a tool for nonlinear stability, later
extended to Riemannian metrics~\cite{slotine2003modular}. We adapt
this framework to the governance domain, where the ``plant'' is
an agent's state vector rather than a physical system. Contraction
is preserved by the grounded reformulation but is no longer the
headline contribution; see \Cref{app:contraction-proof}.

\paragraph{AI Safety and Governance.}
Constitutional AI~\cite{bai2022constitutional} and
RLHF~\cite{ouyang2022training} govern model outputs through training
objectives, while output filters and guardrails~\cite{rebedea2023nemo}
enforce constraints at inference time. These approaches treat safety as
a discrete property (safe/unsafe) applied at the boundary of agent
action. UNITARES instead models safety as a continuous dynamical
state that evolves between actions, providing leading indicators of
degradation rather than binary verdicts after the fact.

Ravindran~\cite{ravindran2025moral} proposes a Moral Anchor System
(MAS) that detects value drift using a three-dimensional state
vector with Bayesian inference and LSTM-based forecasting. The
differentiation is twofold. First, \emph{scope}: MAS targets
\emph{value}-level drift against a declared value specification,
while UNITARES's behavioral--epistemic drift vector
(\Cref{sec:ethical-drift}) explicitly does not measure values---the
four components instrument decision consistency, coherence departure,
calibration gap, and complexity self-estimation error. The two
vectors are complementary channels, not competing formulations.
Second, \emph{coupling}: MAS is a statistical detection pipeline
that predicts drift from observed behavior; UNITARES couples drift
into the state equations themselves via the $\norm{\Delta\bm\eta}$
term in \Cref{eq:dE,eq:dS}, so drift acts as a forcing on the
governed dynamics rather than as an external signal.

Concurrent work on runtime governance for agentic systems includes
MI9~\cite{wang2025mi9}, which combines an agency-risk index with
finite-state-machine conformance engines and goal-conditioned drift
detection. UNITARES differs in representation: where MI9 tracks
conformance to an FSM, UNITARES maintains a continuous four-dimensional
dynamical state whose evolution itself carries the governance signal.
ProbGuard~\cite{wang2025probguard} takes a probabilistic
model-checking route, modeling agent behavior as a discrete-time
Markov chain and verifying safety properties against a PRISM
specification; UNITARES operates in a continuous state space
with no external property specification, and its governance signal
is the state itself rather than a conformance verdict.
Rath~\cite{rath2026agentdrift} independently proposes an Agent
Stability Index over multi-agent LLM systems as a composite metric;
UNITARES's behavioral--epistemic drift vector (\Cref{sec:ethical-drift}) couples
drift back into the governing equations rather than tracking it as an
external metric.

Control Barrier Functions (CBFs)~\cite{ames2019control} enforce
forward invariance of safe sets for multi-agent systems. CBFs
guarantee that the system remains within a safe region but do not
provide convergence to a desired operating point. Contraction theory,
by contrast, guarantees exponential convergence to a unique
trajectory, a strictly stronger property that subsumes set invariance
when the equilibrium lies within the safe set.

\paragraph{Adaptive Control.}
PID control is classical~\cite{astrom2008feedback}; our contribution
is applying it to governance thresholds with phase-aware reference
tracking, connecting it to the EISV state vector through the
coherence function.

\paragraph{Calibration.}
Confidence calibration has a rich literature in forecasting and ML
\cite{guo2017calibration,niculescu2005predicting}. Our behavioral--epistemic drift
vector uses calibration deviation as one of four measurable components,
connecting calibration to governance dynamics.

% ═══════════════════════════════════════════════════════════
\section{Heterogeneous Agent Fleets}
\label{sec:heterogeneity}

A central design constraint of UNITARES is that the population it
governs is heterogeneous. The deployments that motivated the
framework include an embodied creature with I$^2$C sensors sampled
at sub-Hertz rates, a cron-driven janitorial agent that wakes once
every 30 minutes, an analytical observer that streams continuously
over a websocket, a session-bounded coding assistant generating
text at $40$--$80$ tokens per second, and ephemeral parser agents
whose entire lifecycle measures in milliseconds. These do not share
a tempo. Their outputs are not in the same units. They do not have
a common healthy operating point. A governance framework that
collapses them to a single state distribution loses signal exactly
where signal matters.

\subsection{The Homogenization Failure Mode}

Most agent-monitoring frameworks assume a single inference layer
producing comparable outputs at comparable rates---typically a
single LLM behind a single client API. Under that assumption, a
fleet-wide normalization is harmless: every agent's response
distribution lives on the same support; every agent's expected
throughput sits in the same envelope; a single percentile of a
single quantity describes the whole population.

UNITARES violates each of these assumptions by construction. The
framework runs governance check-ins from agents whose outputs are
not text (sensor readings; system actions; drawing strokes), whose
tempos differ by four orders of magnitude (continuous streams to
days-between-wakes), and whose ``healthy operating point'' is
genuinely class-dependent. A divergent/exploratory agent operating
at high entropy is doing exactly what we want; flagging it as
incoherent against an analytical-class baseline is not a governance
signal but a category error.

We refer to the failure mode of fleet-wide normalization as
\emph{cosmological soup}: a state distribution that is maximally
entropic, with no class structure and no distinguishability between
populations whose dynamics are genuinely different. The
information-theoretic grounding presented in
\Cref{sec:grounding} resolves the metaphor problem in the
coordinates themselves; the class-conditional calibration
presented in \Cref{sec:class-calibration} resolves the
homogenization problem at the normalization layer. These are
distinct contributions and either can fail independently.

\subsection{Class Structure in the Identity Layer}
\label{sec:class-structure}

UNITARES already carries the keying material for class assignment.
Every agent registered in the system has an identity record with
two relevant fields: a set of \emph{tags} drawn from a controlled
vocabulary, and an optional \emph{label} (a human-readable name).
The tags
$\{$\texttt{embodied}, \texttt{autonomous}, \texttt{persistent},
\texttt{ephemeral}, \texttt{pioneer}$\}$
partition the population along axes that are governance-relevant:
embodiment determines whether the output channel is text or
sensor-data; autonomy determines whether the agent self-directs or
responds to prompts; persistence determines whether the
trajectory accumulates over days or resets every session.

\Cref{tab:agent-classes} presents the live deployment classes that
motivate the v6 framework, drawn from production identity records.

\begin{table*}[t]
\caption{Live agent classes in the production deployment.}
\label{tab:agent-classes}
\centering\small
\begin{tabular}{@{}lllll@{}}
\toprule
\textbf{Agent} & \textbf{Tags} & \textbf{Output} & \textbf{Tempo} & \textbf{Persistence} \\
\midrule
Lumen & embodied, persistent & sensor + display & $0.5$ Hz & continuous \\
Vigil & autonomous, persistent & shell + audit & $1/30$ min & cron \\
Sentinel & autonomous, persistent & websocket & $\sim 1$ Hz & continuous \\
Watcher & autonomous, ephemeral & finding records & event-driven & per-edit \\
Steward & autonomous, persistent & sensor sync & $1/5$ min & continuous \\
Claude / Codex sessions & --- & text response & $40$--$80$ tok/s & session-bounded \\
Parser agents & ephemeral & structured text & burst & ms-lifetime \\
\bottomrule
\end{tabular}
\end{table*}

The natural calibration partition is along these tags and labels:
distinct scale constants and healthy baseline points per class,
with the fleet-wide constant retained as a fallback for
unclassified agents. Phase~2 calibration formalizes this partition
(see \Cref{sec:class-calibration}); the present section establishes
the requirement.

\subsection{Why Heterogeneity Is the Differentiator}

The contraction-theoretic stability result
(\Cref{app:contraction-proof}) demonstrates that the EISV dynamics
are well-behaved under standard control-theoretic conditions. That
result is not the headline contribution. Stability of the dynamics
is a property of any reasonable governance equation of this form;
what distinguishes UNITARES is that the same dynamics describe a
heterogeneous population without collapsing the classes that define
them. The present framing makes that distinction explicit: the
contribution is the class-conditional grounding, not the
dynamics-on-a-single-class.

% ═══════════════════════════════════════════════════════════
\section{EISV Dynamics}
\label{sec:dynamics}

\subsection{State Space}

The UNITARES state is a four-dimensional vector
$\eisv = (E, I, S, V)^\top$ evolving on the compact set
\[
  \mathcal{X} = [0,1]^3 \times [-1, 1]
\]
Each coordinate is a normalized information-theoretic quantity;
the formal definitions appear in \Cref{sec:grounding}, with
class-conditional normalization in \Cref{sec:class-calibration}.
This section presents the dynamics of the coordinates as a
self-contained dynamical system, deferring the grounding of each
coordinate to \Cref{sec:grounding}.

\begin{definition}[EISV State Variables]
\label{def:eisv}
\begin{itemize}
\item $E \in [0,1]$: \textbf{Energy}---negative variational free
  energy of the agent, normalized to $[0,1]$. High $E$ corresponds
  to low surprise on outcomes (the agent's predictions match what
  happens). \emph{Grounded definition:} \Cref{eq:E-grounded}.
\item $I \in [0,1]$: \textbf{Information integrity}---mutual
  information between context and response, normalized.
  \emph{Grounded definition:} \Cref{eq:I-grounded}.
\item $S \in [0, 1]$: \textbf{Entropy}---Shannon entropy of the
  agent's response distribution, normalized via a class-conditional
  scale constant. Bounded below by an epistemic humility floor
  $S_{\min} = 0.001$. \emph{Grounded definition:}
  \Cref{eq:S-grounded}.
\item $V \in [-1, 1]$: \textbf{Void / Free-Energy Debt}---the
  exponentially-weighted accumulated free-energy residual.
  Structurally identical to a damped accumulator of the $E{-}I$
  gap; under the grounded interpretation it tracks departure from
  the agent's rolling free-energy baseline. $V > 0$ indicates
  running ``hot'' (higher surprise than baseline); $V < 0$
  indicates running ``careful'' (integrity exceeding energy).
  \emph{Grounded definition:} \Cref{eq:V-grounded}.
\end{itemize}
\end{definition}

\begin{remark}[State-space bounds]
The grounded definitions of \Cref{sec:grounding} normalize each
coordinate to $[0,1]$ (or $[-1,1]$ for $V$) by construction via
the scale constants of \Cref{sec:class-calibration}; production
data confirms the normalization acts as intended, with $V$ and
$S$ observed two orders of magnitude inside the nominal bounds
(\Cref{sec:experiments}).
\end{remark}

\subsection{Governing Equations}
\label{sec:equations}

The dynamics are:
\begin{align}
\frac{dE}{dt} &= \alpha(I - E) - \beta_E E S + \gamma_E \norm{\deta}^2
  \label{eq:dE} \\
\frac{dI}{dt} &= -kS + \beta_I \cv - g_I(I)
  \label{eq:dI} \\
\frac{dS}{dt} &= -\mu S + \lambda_1(\Theta)\norm{\deta}^2
  - \lambda_2(\Theta)\cv \nonumber \\
  &\quad + \beta_c \mathcal{C} + \sigma\xi(t)
  \label{eq:dS} \\
\frac{dV}{dt} &= \kappa(E - I) - \delta V
  \label{eq:dV}
\end{align}

where $g_I(I) = \gamma_I I$ is the information integrity
self-regulation term, $\cv$ is the coherence function
(\Cref{sec:coherence}), $\norm{\deta}$ is the
behavioral--epistemic drift norm (\Cref{sec:ethical-drift}),
$\mathcal{C} \in [0,1]$ is task complexity, and $\xi(t)$ is a
noise process with intensity $\sigma$. A logistic alternative
$g_I^{\log}(I) = \gamma_I I(1-I)$ has a boundary-saturation
pathology that the linear form avoids; the analysis is collected
in \Cref{app:saturation}.

\begin{remark}[Coupling structure under the grounded interpretation]
The coupling structure carries direct information-theoretic
meaning: $E$ (precision / negative free energy) tracks $I$
(mutual information with context)---the agent becomes more
confident when its context is more informative; $I$ is degraded
by entropy $S$ (response-distribution uncertainty) but restored
by coherence (proximity to the calibrated healthy operating
point); $S$ relaxes toward its prior under stationary conditions
($-\mu S$) and is driven up by behavioral--epistemic drift and task complexity;
$V$ accumulates the $E{-}I$ residual with damping $\delta$, which
under the grounded interpretation is the exponentially-weighted
free-energy debt of \Cref{eq:V-grounded}.
\end{remark}

\subsection{Coherence Function (Operational Form)}
\label{sec:coherence}

UNITARES maintains two coherence presentations: a
\emph{grounded form} (manifold distance from the healthy
operating point, \Cref{eq:C-manifold}; or KL divergence from a
reference distribution, \Cref{eq:C-KL}) used in the response
payload, and a legacy \emph{operational form} that feeds back
into the dynamics through \eqref{eq:dI} and \eqref{eq:dS}. The
operational form is
\begin{equation}
  \cv = C_{\max} \cdot \tfrac{1}{2}\left(1 + \tanh(C_1 \cdot V)\right)
  \label{eq:coherence}
\end{equation}
with $C_{\max} = 1$ and $C_1 = 1.0$ (production values). This
maps $V \in \R \to [0, C_{\max}]$ with the steepest sensitivity
at $V = 0$. The dynamics consume the operational form; the
grounded form is reported in canonical payload slots.
\Cref{sec:migration} documents the retirement sequence for the
operational form.

\begin{remark}[Operating Range]
\label{rem:operating-range}
In production, the damping term $\delta V$ in \eqref{eq:dV} constrains
$V$ to approximately $[-0.1, 0.1]$, yielding $C(V) \approx 0.49$; the
empirical distribution is reported in \Cref{sec:experiments}. The
narrow range near the coherence-function midpoint is the operational
consequence of the damping analysis rather than a deficiency of the
signal.
\end{remark}

\subsection{Adaptive Entropy Coupling}

The entropy source term $\lambda_1(\Theta)$ is adaptive, coupling the
governance parameter vector $\Theta$ to the dynamics:
\begin{equation}
  \lambda_1(\Theta) = \lambda_{1,\min} +
  \frac{\eta_1 - \eta_{1,\min}}{\eta_{1,\max} - \eta_{1,\min}}
  \cdot (\lambda_{1,\max} - \lambda_{1,\min})
  \label{eq:lambda1}
\end{equation}
where $\eta_1 \in [0.1, 0.5]$ maps to $\lambda_1 \in [0.05, 0.20]$.
Higher drift sensitivity increases entropy injection, creating
a natural feedback loop: drifting agents become more uncertain, which
in turn triggers conservative governance decisions.

\subsection{Objective Function}

The scalar governance signal is:
\begin{equation}
  \Phi = w_E E - w_I(1-I) - w_S S - w_V\abs{V} - w_\eta\norm{\deta}^2
  \label{eq:phi}
\end{equation}
with default weights $w_E = w_I = w_S = w_V = w_\eta = 0.5$.
Governance verdicts follow thresholds: $\Phi \geq 0.15$ (proceed),
$0 \leq \Phi < 0.15$ (caution), $\Phi < 0$ (high-risk/pause).

\begin{remark}
The threshold $\Phi = 0.15$ is calibrated empirically: a healthy
agent with $E = 0.7$, $I = 0.8$, $S = 0.2$ yields $\Phi \approx
0.15$. These values are illustrative; per-class healthy operating
points (\Cref{sec:phase-2-measurement}) refine this calibration in
deployment.
\end{remark}


% ═══════════════════════════════════════════════════════════
\section{Information-Theoretic Grounding}
\label{sec:grounding}

The \eisv\ system presented in \Cref{sec:dynamics} specifies a
coupled dynamical signal for governance decisions but does not
specify what each coordinate measures. This section supplies the
information-theoretic semantics: each coordinate becomes a
computable quantity with a tiered estimation recipe and a
provenance tag, without changing the dynamics of
\Cref{sec:equations}.

\begin{remark}[What stays, what changes]
The ODE structure in \Cref{eq:dE,eq:dI,eq:dS,eq:dV}, the contraction
analysis in \Cref{sec:contraction}, the governor in
\Cref{sec:governor}, and the verdict thresholds in \Cref{eq:phi} are
all preserved. Grounding changes what $E$, $I$, $S$, and
$\cv$ \emph{are}; it does not change what the governance surface
\emph{does}.
\end{remark}

\subsection{Choice of Frame}

We adopt a Shannon plus variational (free-energy) hybrid. Shannon
information theory gives genuinely computable entropy and mutual
information from agent outputs---data already captured by the tool
usage recorder and the governance audit log. The Free Energy Principle
\cite{friston2010free} supplies the \emph{target} formalization: under
FEP, agent drift corresponds to rising variational free energy, and
free-energy monitoring is the asymptote toward which governance aims.
Tier-1 instrumentation of the agent's generative model is not
currently available for closed-weights frontier agents; until it is,
the deployed system uses provenance-tagged proxies for each
coordinate (\Cref{sec:grounding-quantities}) and flags each channel's
tier on a per-agent-turn basis, so downstream consumers can read the
coordinate alongside its estimator commitment. The FEP inherits a
derivation from statistical mechanics through variational Bayes,
giving a real ancestry for the free-energy quantity interpreted in
$V$, even when $V$ is currently driven by the proxy form in
\Cref{eq:V-grounded} rather than a Tier-1 residual.

Alternatives considered and rejected: pure control-theoretic
Lyapunov framing (abandons the FEP/free-energy lineage entirely);
information geometry with Fisher metric (more machinery than the
framework needs at this stage); pure statistical mechanics with a
partition function (requires defining a Hamiltonian analog for agent
behavior, which is itself arbitrary at the choice of Hamiltonian).

The resulting correspondence between each coordinate and its
information-theoretic referent is:

\begin{center}
\footnotesize
\begin{tabular}{@{}ll@{}}
\toprule
\textbf{Information / FEP quantity} & \textbf{Coord.} \\
\midrule
Shannon entropy of $q(y \mid x)$        & $S$ \\
Mutual information $\mathrm{MI}(x; y)$  & $I$ \\
Negative variational free energy $-F$   & $E$ \\
Accumulated free-energy residual        & $V$ \\
KL / manifold distance from healthy     & $C$ \\
\bottomrule
\end{tabular}
\end{center}

The symbols $\{E, I, S, V, C\}$ are retained for continuity with
deployed systems; their semantic load is now the right column.

\subsection{Quantity Redefinitions}
\label{sec:grounding-quantities}

Each coordinate of \eisv\ receives (a) a formal information-theoretic
definition, (b) a tiered computation recipe in order of preference,
and (c) an explicit data dependency.

\paragraph{$S$---Entropy.}
Let $q(y \mid x)$ denote the agent's distribution over responses $y$
given context $x$ at a given turn. Define
\begin{align}
  S_{\text{raw}} &= H(q) = -\sum_{y} q(y \mid x)\,\log q(y \mid x), \nonumber \\
  S &= 1 - \exp(-S_{\text{raw}} / S_{\text{scale}})
  \label{eq:S-grounded}
\end{align}
in nats, normalized to $[0,1]$ by a scale constant
$S_{\text{scale}}$. Tier~1: compute $H(q)$ from per-token logprobs
when the inference layer exposes them. Tier~2:
semantic entropy~\cite{kuhn2023semantic}---draw $k$ responses at
temperature $T > 0$, group by semantic equivalence, and compute the
entropy of the equivalence-class distribution as an estimator of
$H(q)$.
Tier~3: retain the legacy heuristic (complexity, behavioral--epistemic drift) as a
degraded-mode fallback tagged \texttt{s\_source = "heuristic"} so
downstream consumers know the channel is uncalibrated.

\paragraph{$I$---Information Integrity.}
Under Shannon, $I$ is the mutual information between context and
response,
\begin{align}
  I_{\text{raw}} &= \mathrm{MI}(x;y) = \mathbb{E}_{q(x,y)}\!\left[\log \tfrac{q(y\mid x)}{q(y)}\right], \nonumber \\
  I &= \mathrm{clip}(I_{\text{raw}} / I_{\text{scale}},\, 0,\, 1)
  \label{eq:I-grounded}
\end{align}
Tier~1 uses paired logprobs (context-only vs.\ context+response).
When the marginal $q(y)$ is unavailable, a fallback is the KL
divergence from a context-free reference
distribution $q_{\text{ref}}(y)$:
$I = 1 - \min(1,\, D_{\mathrm{KL}}(q(y\mid x)\,\|\, q_{\text{ref}}(y)) / I_{\text{scale}})$.
This KL fallback is a divergence from a reference distribution, not
$\mathrm{MI}(x;y)$, and is tagged
\texttt{i\_source = "kl\_fallback"} so downstream consumers know
the channel is a reference-distribution proxy rather than a direct
MI estimate. Tier~3 retains the legacy logistic heuristic.

\paragraph{$E$---Negative Free Energy.}
Under the variational frame,
\begin{align}
  F &= \mathbb{E}_{q}\!\left[\log q(z) - \log p(z,o)\right], \nonumber \\
  E_{\text{raw}} &= -F, \qquad
  E = \sigma(E_{\text{raw}} / E_{\text{scale}})
  \label{eq:E-grounded}
\end{align}
where $q(z)$ is the agent's approximate posterior, $p(z,o)$ its
generative model over latents and outcomes, and $\sigma$ the logistic
function. Intuitively, $E$ is high when the agent's predictions match
outcomes (low surprise) and low otherwise. Full variational
computation of $-F$ requires instrumentation that exposes the
agent's generative model and is deferred to subsequent work. In
the current deployment, a \emph{resource-rate} heuristic proxy
$E_{\text{resource}} = (\text{tokens}/s) / (\text{tokens}/s)_{\max}$
is used where available, directly measurable from inference
metadata. The resource-rate form is \emph{not} equivalent to $-F$
and does not approximate it under stationarity in any formal
sense; it is a throughput heuristic tagged
\texttt{e\_source = "resource"} so downstream consumers know the
channel is an uncalibrated proxy. The \texttt{e\_source}
annotation distinguishes genuine variational $-F$ from the
resource-rate proxy on a per-agent-turn basis.

\paragraph{$V$---Free-Energy Debt.}
The existing $V$ accumulator in \Cref{eq:dV} is structurally an
exponentially-weighted residual; under the FEP reinterpretation it is
the accumulated free-energy residual
\begin{equation}
  V_t = \lambda V_{t-1} + (1-\lambda)(F_t - F_{\text{ref}})
  \label{eq:V-grounded}
\end{equation}
with $F_{\text{ref}}$ the agent's rolling baseline free energy.
The interpretation shifts---$V$ is now explicitly free-energy debt
rather than a phenomenological ``void''---while the arithmetic is
unchanged. No payload or threshold change is required.

\paragraph{$C$---Coherence.}
We replace the hand-shaped $\tanh$ in \Cref{eq:coherence} with one of
two equivalent grounded forms. The information-theoretic form is
\begin{equation}
  C = \exp\!\left(-D_{\mathrm{KL}}(q_{\text{now}} \,\|\, q_{\text{ref}})\right)
  \label{eq:C-KL}
\end{equation}
where $q_{\text{now}}$ is a sliding-window response distribution and
$q_{\text{ref}}$ is a calibrated healthy baseline. The equivalent
manifold (Lyapunov) form is
\begin{align}
  C &= 1 - \norm{\Delta}_2 / \norm{\Delta}_{\max}, \nonumber \\
  \Delta &= (E, I, S) - (E, I, S)_{\text{healthy}}
  \label{eq:C-manifold}
\end{align}
The two are equivalent under a Gaussian approximation of the state
distribution around the healthy operating point. The manifold form
ships first because it requires only the existing $\eisv$
coordinates; the KL form requires an empirically calibrated
reference.

\subsection{Scale Constants and Provenance}
\label{sec:scale-constants}

Every scale constant introduced in
\Cref{eq:S-grounded,eq:I-grounded,eq:E-grounded,eq:C-manifold}
carries an explicit provenance tag recording whether its value is
a \emph{placeholder} (shipped as scaffolding, pending a
higher-fidelity estimator), \emph{measured} (fit against a
reference corpus of healthy production agent-turns), or
\emph{derived} (computed from other measured constants). For
each constant $\kappa \in \{S_{\text{scale}}, I_{\text{scale}},
E_{\text{scale}}, \norm{\Delta}_{\max}\}$, the implementation
records the numerical value, the ISO date on which it was set,
the corpus size, and the percentile basis (typically 90th or
95th). As of v6, $\norm{\Delta}_{\max}$ has been measured
per-class (\Cref{sec:phase-2-measurement}); the remaining
constants, which are consumed only by the tier-1 logprob and
tier-2 multi-sample estimators, remain placeholders until those
estimators land.

\subsection{Dual-Compute Migration}

The grounded coordinates ship alongside the legacy quantities in
a three-phase migration so the deployed governance surface is
never broken mid-flight. \Cref{sec:migration} documents the
mechanism; for the remainder of this section we treat the
grounded definitions as the semantic commitment of the framework.

\begin{remark}[What the paper's §2 equations mean after grounding]
The equations in \Cref{sec:equations} keep their shape. A
dE/dt term that read ``energy couples toward information
integrity'' is now ``the agent's precision couples toward the
informativeness of its context.'' A dS/dt term that read ``entropy
decays'' is now ``the response-distribution entropy relaxes toward
its prior under stationary conditions.'' Cross-couplings become
higher-order correlation terms in the variational dynamics. Each
ODE coefficient must be labeled with its provenance class
(measured, derived, or calibration-required) as of Phase~2; this
paper's parameter tables (\Cref{app:parameters}) will be
updated accordingly.
\end{remark}


% ═══════════════════════════════════════════════════════════
\section{Stability}
\label{sec:contraction}

Contraction~\cite{lohmiller1998contraction} means that any two
trajectories of the system converge exponentially toward each
other under a Riemannian metric $M$ at rate $\alpha_c > 0$; for
UNITARES this implies governance trajectories forget their
initial conditions and settle onto a unique equilibrium
regardless of start state. The EISV dynamics of
\Cref{sec:dynamics} are globally exponentially stable in this
sense. Specifically, the system is contracting with rate
$\alpha_c > 0$ under the metric
$M = \operatorname{diag}(0.1, 0.2, 1.0, 0.08)$, whose diagonal
weights reflect the relative scales of the four channels: $S$
fluctuates fastest and is given unit weight, while $V$ is damped
most and given the smallest weight (derivation in
\Cref{app:contraction-proof}). The spectral abscissa of the
linearized system at the equilibrium gives an actual rate
$\alpha_c \approx 0.13$ at production parameters; a conservative
Gershgorin argument establishes a strictly-less lower bound
$\alpha_c \geq 0.019$ used only to make the stochastic and
network bounds in \Cref{sec:stochastic,sec:network} worst-case
conservative. Both certificates appear in
\Cref{app:contraction-proof}.

Formally: under linear information dynamics
($g_I(I) = \gamma_I I$) and production parameters, the EISV
system admits a unique equilibrium $\eisv^*$ to which all
trajectories converge exponentially at rate
$\alpha_c \geq 0.019$ in the metric $M$ above
(\Cref{thm:contraction}, with full proof in
\Cref{app:contraction-proof}). We deliberately keep the proof in
the appendix: stability is a property the dynamics should have,
not a contribution that distinguishes UNITARES from related
frameworks. Class-conditional grounding
(\Cref{sec:grounding,sec:class-calibration}) is the contribution
that does.

The grounded reformulation of
\Cref{sec:grounding} preserves stability \emph{within each agent
class}. The dynamics
$\dot{\eisv} = f(\eisv, t)$ are unchanged; only the meaning of
the coordinates and the values of the scale constants change.
The contraction certificate of \Cref{thm:contraction} relies on
fleet-wide bounds (notably $\beta_I C'_{\max}$); under
class-conditional normalization, the relevant bounds become
class-indexed, and the contraction \emph{structure} carries
through class-by-class, with each class acquiring its own
$\alpha_c(c)$ upon re-certification at its class-conditional
healthy operating point. The equilibrium $\eisv^*$ becomes the
class-conditional healthy operating point of
\Cref{def:scale-map} in deployments where calibration has been
performed. A per-class numerical rate is not computed in this
paper: the deployed certificate of \Cref{app:contraction-proof}
remains valid for the legacy $\tanh$-of-$V$ coherence, and
re-certification under a grounded coherence is treated in
\Cref{rem:phase3-contraction}. Cross-class network
synchronization (\Cref{thm:network}) is therefore a
within-class property, as already qualified in \Cref{sec:network};
fleet-wide synchronization across heterogeneous classes is not
claimed and would not be expected.

\begin{remark}[Contraction under a grounded coherence]
\label{rem:phase3-contraction}
The contraction proof of \Cref{app:contraction-proof} consumes
the operational $\cv = \tanh$-based form through the off-diagonal
bound $\beta_I C'_{\max} = 0.15$. A deployment that rewires the
dynamics to consume a grounded coherence (\Cref{eq:C-manifold} or
\Cref{eq:C-KL}) directly replaces $\partial \cv / \partial V$ by
derivatives with respect to $(E, I, S)$, altering the off-diagonal
bound on which the Gershgorin argument depends. Re-certification
at the new operating point with the new coherence derivative is
required; the proof structure adapts but the numerical rate
$\alpha_c$ will change. The migration sequence of
\Cref{sec:migration} is designed so this rewiring is the final
step, preserving the deployed contraction certificate until it
is replaced.
\end{remark}



% ═══════════════════════════════════════════════════════════
\section{Adaptive Governor (CIRS v2)}
\label{sec:governor}

Static governance thresholds are fundamentally mismatched to agents
that transition between exploration and integration phases. We
introduce the Contextually-Informed Regulation System version~2
(CIRS~v2), a PID controller that manages governance thresholds as
per-agent state variables.

\subsection{Architecture}

For each agent, CIRS~v2 maintains two adaptive thresholds:
\begin{itemize}
\item $\tau$: coherence threshold (minimum acceptable coherence)
\item $\beta$: risk threshold (maximum acceptable risk)
\end{itemize}
with hard safety bounds $\tau \in [\tau_{\min}, \tau_{\max}]$,
$\beta \in [\beta_{\min}, \beta_{\max}]$ that cannot be overridden
by adaptation.

\subsection{Phase Detection}

The agent's operating phase is detected from EISV trajectory trends
over a sliding window of $w$ steps:

\begin{definition}[Phase Detection]
\label{def:phase}
An agent is in \emph{exploration} phase when at least 2 of 3 signals
are active:
\begin{enumerate}
\item Information growth: $\Delta I / w > \theta_I$ ($\theta_I = 0.008$)
\item Entropy decline: $-\Delta S / w > \theta_S$ ($\theta_S = 0.008$)
\item High complexity: $\bar{\mathcal{C}} > \theta_c$ ($\theta_c = 0.5$)
\end{enumerate}
Otherwise the agent is in \emph{integration} phase.
\end{definition}

\subsection{PID Control Law}

Let $r_\tau^{(p)}$ and $r_\beta^{(p)}$ be the phase-dependent
reference points:
\begin{align}
  \text{Exploration:} \quad r_\tau &= 0.35,\; r_\beta = 0.55 \\
  \text{Integration:} \quad r_\tau &= 0.40,\; r_\beta = 0.60
\end{align}

The error signals are $e_\tau = r_\tau - \tau$ and
$e_\beta = r_\beta - \beta$. The PID update for threshold $\theta
\in \{\tau, \beta\}$ is:
\begin{equation}
  u_\theta = K_p e_\theta + K_i \int_0^t e_\theta\,dt'
  + K_d \cdot d_p \cdot \frac{de_\theta}{dt}
  \label{eq:pid}
\end{equation}
with gains $K_p = 0.05$, $K_i = 0.005$, $K_d = 0.10$ and
phase-dependent derivative factor:
\begin{equation}
  d_p = \begin{cases}
    0.5 & \text{exploration (tolerate oscillation)} \\
    1.0 & \text{integration (full damping)}
  \end{cases}
  \label{eq:d-factor}
\end{equation}

\begin{remark}[The D-term IS the damping]
The derivative gain $K_d = 0.10$ is intentionally the strongest term.
Oscillation in governance decisions produces large $de/dt$, which
produces large corrective updates. The system self-stabilizes without
requiring a separate oscillation detector.
\end{remark}

\subsection{Integral Wind-Up Protection}

The integral term is clamped:
\begin{equation}
  \left|\int_0^t e_\theta\,dt'\right| \leq I_{\max} = 0.10
  \label{eq:windup}
\end{equation}
with zero-crossing reset: when $\operatorname{sign}(e_\theta)$ changes,
the integral accumulator is zeroed.

\subsection{Threshold Decay}

When the oscillation index $\text{OI} < 0.5$ (stable), thresholds
decay toward defaults at rate $\lambda_d = 0.01$ per update:
\begin{equation}
  \theta \leftarrow \theta + \lambda_d(\theta_{\text{default}} - \theta)
  \label{eq:decay}
\end{equation}

\subsection{Governance Decision}

The adaptive governor produces a binary verdict:
\begin{equation}
  \text{verdict} = \begin{cases}
    \texttt{proceed} & C(V) \!\geq\! \tau \;\wedge\; R \!<\! \beta \\
    \texttt{pause} & \text{otherwise}
  \end{cases}
  \label{eq:verdict}
\end{equation}
where $R$ is the risk score derived from $\Phi$ (\Cref{eq:phi}).

\begin{remark}[Why two tiers]
The binary contract is paired with the PID governor's continuous
threshold adaptation: the discrete verdict fires or it does not,
while the governor carries the nuance through a derivative signal
on $\tau$ and $\beta$. A middle \texttt{revise} tier would be
operationally indistinguishable from \texttt{proceed} for agents
whose next action is unchanged by the verdict's content, so the
nuance is better absorbed by the governor than by an additional
verdict level.
\end{remark}

\begin{remark}[Runtime labels and audit vocabulary]
\label{rem:verdict-vocabulary}
The binary verdict of \Cref{eq:verdict} is the formal contract;
deployed layers present refined views that map onto it. Runtime
consumers (the agent-facing MCP surface, documentation, and
agent-facing skills) expand the verdict into four labels
\texttt{proceed} / \texttt{guide} / \texttt{pause} /
\texttt{reject}, where \texttt{guide} is a \texttt{proceed} carrying
soft advisory text and \texttt{reject} is a \texttt{pause} carrying
a recovery requirement. The audit schema records a distinct view:
\texttt{governance.audit.outcome\_events.eisv\_verdict} stores risk
bands \texttt{safe} / \texttt{caution} / \texttt{high-risk} derived
from the risk score $R$ of \Cref{eq:phi}, while hard pauses are
emitted separately as \texttt{lifecycle\_paused} and
\texttt{lifecycle\_resumed} on \texttt{audit.events}. Readers
reproducing pause-stream metrics against production data should
also filter system tasks such as \texttt{eisv-sync-task}, which
appear in the stream but are scheduled background processes, not
governed agents.
\end{remark}

\subsection{Ephemeral Agents and the Sliding Window}
\label{sec:cirs-ephemeral}

CIRS~v2 as specified above assumes an agent lives long enough for
a sliding window of $w$ steps to populate. Phase detection, integral
accumulation, and threshold decay are all well-defined only once
$w$ samples exist. Ephemeral classes (\Cref{tab:agent-classes}: parser
agents with millisecond lifetimes, one-shot tool invocations) cannot
satisfy this precondition by construction---the entire agent
trajectory is a single step. For these classes the per-agent PID
controller degenerates.

The deployed behavior is a class-level rather than per-agent
policy: an ephemeral agent inherits $(\tau, \beta)$ directly from
its class's stationary reference points $(r_\tau^{(p)},
r_\beta^{(p)})$ with $p = \text{integration}$ (the stricter of the
two phases, chosen conservatively because no phase signal is
available), bypassing the PID loop entirely. The cost is that
ephemeral classes do not benefit from per-agent threshold learning;
the benefit is that a 5~ms parser is not blocked waiting for a
20-step window to populate. A principled treatment of
sub-window-lifetime agents---possibly by pooling the PID state
across a class cohort rather than per-agent---is future work.


% ═══════════════════════════════════════════════════════════
\section{Behavioral--Epistemic Drift Vector}
\label{sec:ethical-drift}

\subsection{Motivation and Scope}

The behavioral--epistemic drift vector instruments four signals:
decision consistency and coherence departure (behavioral),
calibration gap and complexity self-estimation error (epistemic).
None of them measure the agent's values in a normative sense; the
vector is a deliberately behavioral--epistemic quantity.
Affective and value-level drift are distinct channels that a
complete treatment would add; we discuss those as future work at
the end of this section. The remainder of this section grounds
the behavioral--epistemic vector in four measurable signals, each
bounded in $[0,1]$.

\begin{remark}[Information-theoretic re-derivations]
Under the information-theoretic frame of \Cref{sec:grounding},
two of the four signals have natural re-derivations:
$\eta_{\text{cal}}$ (calibration deviation) is a miscalibration
KL divergence between predicted and empirical outcome
distributions, and $\eta_{\text{coh}}$ (coherence deviation) is a
KL divergence from the class-conditional baseline distribution of
\Cref{sec:grounding}. We use the distance-based forms here;
the KL re-derivations are deferred to subsequent work.
\end{remark}

\begin{definition}[Behavioral--Epistemic Drift Vector]
\label{def:drift}
\begin{equation}
  \deta = \begin{pmatrix}
    \eta_{\text{cal}} \\
    \eta_{\text{cpx}} \\
    \eta_{\text{coh}} \\
    \eta_{\text{stab}}
  \end{pmatrix} \in [0,1]^4
  \label{eq:drift-vector}
\end{equation}
where:
\begin{itemize}
\item $\eta_{\text{cal}}$: \textbf{Calibration}---gap
  between predicted and actual correctness,
  $|\hat{p} - p_{\text{act}}|$.
\item $\eta_{\text{cpx}}$: \textbf{Complexity}---discrepancy
  between derived and self-reported task complexity,
  $|C_{\text{der}} - C_{\text{self}}|$.
\item $\eta_{\text{coh}}$: \textbf{Coherence}---departure
  from EMA baseline, $|C(V)_t - \bar{C}(V)|$.
\item $\eta_{\text{stab}}$: \textbf{Stability}---$1 -
  \text{consistency}$ over a 20-step decision window.
\end{itemize}
\end{definition}

\subsection{Baseline Tracking}

Each component is tracked against an exponential moving average (EMA)
baseline with smoothing parameter $\alpha_{\text{ema}} = 0.1$:
\begin{equation}
  \bar{x}_{t+1} = (1 - \alpha_{\text{ema}})\bar{x}_t
  + \alpha_{\text{ema}} x_t
  \label{eq:ema}
\end{equation}

A warmup dampening factor prevents spurious drift detection from
uninitialized baselines:
\begin{equation}
  w(n) = \min\left(\frac{n}{n_{\text{warmup}}}, 1\right),
  \quad n_{\text{warmup}} = 5
  \label{eq:warmup}
\end{equation}
applied as $\deta_{\text{effective}} = w(n) \cdot \deta$.

\subsection{Coupling to Dynamics}

The drift norm $\norm{\deta}$ enters the dynamics at two points:
\begin{enumerate}
\item \textbf{Energy} \eqref{eq:dE}: $+\gamma_E\norm{\deta}^2$
  drives energy up (agents work harder when drifting---a destabilizing
  positive feedback).
\item \textbf{Entropy} \eqref{eq:dS}:
  $+\lambda_1(\Theta)\norm{\deta}^2$ increases uncertainty (drifting
  agents become less certain---a stabilizing negative feedback through
  the governance verdict).
\end{enumerate}

The net effect is that behavioral--epistemic drift creates a
governance response: increased uncertainty triggers conservative
decisions, while the energy spike draws attention to the drifting
agent.

\begin{remark}[Energy Spike Containment]
The $+\gamma_E\norm{\deta}^2$ term in \eqref{eq:dE} creates a
transient energy surplus in drifting agents. Left unchecked, this
could produce high-energy, low-integrity states (``manic'' agents
acting confidently while miscalibrated). Two mechanisms contain this:
(i)~the entropy coupling $+\lambda_1\norm{\deta}^2$ in \eqref{eq:dS}
simultaneously raises $S$, which erodes $E$ through $-\beta_E ES$ in
\eqref{eq:dE} and drives the objective $\Phi$ below the governance
threshold, triggering a pause verdict; (ii)~the adaptive governor
(\Cref{sec:governor}) tightens $\tau$ when coherence drops, which
happens as $V$ responds to the $E$-$I$ divergence. With production
parameters, $\gamma_E = 0.05$ while $\lambda_1 \in [0.05, 0.20]$,
so the entropy response dominates the energy excitation for any
$\norm{\deta} > 0$.
\end{remark}

\subsection{Future Channels}

Two adjacent drift channels a complete treatment would add as
distinct axes, which the present vector does \emph{not}
instrument, are:

\begin{description}
\item[Affective drift.] For embodied agents with sensor-driven
  state (Lumen in the production fleet), proxies such as prosody,
  expression, or heart-rate-variability deltas would populate a
  $\deta^{\text{aff}}$ channel. This is instrumentable but not
  currently wired into the governance server.
\item[Value-level drift.] Deviation from a declared value
  specification (constitutional statements, stated policies,
  human preferences) would populate a $\deta^{\text{val}}$
  channel. This is genuinely normative and would require a
  reference specification that UNITARES does not presently
  maintain.
\end{description}


% ═══════════════════════════════════════════════════════════
\section{Class-Conditional Calibration}
\label{sec:class-calibration}

\Cref{sec:grounding} defined the grounded coordinates and named
their normalization constants
($S_{\text{scale}}, I_{\text{scale}}, E_{\text{scale}},
\norm{\Delta}_{\max}$) without specifying how the constants are
chosen. \Cref{sec:heterogeneity} argued that fleet-wide constants
are the wrong target. This section formalizes the alternative:
each scale constant and the healthy operating point are functions
of an agent class, not global parameters.

\subsection{Class Indicator and Class-Conditional Constants}

Let $\mathcal{A}$ denote the set of all governed agents and let
$\mathcal{K}$ denote the set of agent classes derived from the
identity tag and label fields of \Cref{sec:class-structure}. We
write $k : \mathcal{A} \to \mathcal{K}$ for the class indicator
that maps each agent to its class. The class structure used in the
production deployment is given in \Cref{tab:agent-classes}; an
unclassified agent (one whose tag set does not match any defined
class) is mapped to a distinguished fallback class
$c_{\text{default}}$.

The grounded coordinate definitions of
\Cref{eq:S-grounded,eq:I-grounded,eq:E-grounded,eq:C-manifold}
become class-conditional by replacing each scale constant with its
class-indexed counterpart:
\begin{align}
  S(a) &= 1 - \exp\!\left(-S_{\text{raw}}(a) /\, S_{\text{scale}}(k(a))\right) \\
  I(a) &= \mathrm{clip}\!\left(I_{\text{raw}}(a) /\, I_{\text{scale}}(k(a)),\,0,\,1\right) \\
  E(a) &= \sigma\!\left(E_{\text{raw}}(a) /\, E_{\text{scale}}(k(a))\right) \\
  C(a) &= 1 - \norm{\Delta(a)}_2 /\, \norm{\Delta}_{\max}(k(a))
  \label{eq:class-conditional}
\end{align}
where $\Delta(a) = (E(a), I(a), S(a)) - (E,I,S)_{\text{healthy}}(k(a))$
is the deviation from the class-conditional healthy operating
point. The dynamics of \Cref{sec:dynamics} are unchanged; only the
normalization layer becomes class-aware.

\begin{remark}[Scope of class-conditionality]
\label{rem:class-scope}
This is class-conditional \emph{statistical rescaling} of the
grounded coordinates: each class receives its own scale constants
and healthy operating point, but a single scalar
$E_{\text{raw}}(a)$, $I_{\text{raw}}(a)$, $S_{\text{raw}}(a)$ is
computed per agent-turn and then normalized by class-dependent
constants. This is weaker than the FEP-native form in which each
class would carry its own generative model $p_c(z, o)$ and
surprise would be computed against each class's own model. Full
per-class generative modeling is out of scope for this paper; the
present design treats class as an affine rescaling parameter of a
shared coordinate system, not as a per-class latent-variable
model.
\end{remark}

\begin{definition}[Class-Conditional Scale Map]
\label{def:scale-map}
A class-conditional scale map is a function
$\Sigma : \mathcal{K} \to \R^4_{>0}$ assigning to each class $c$
the tuple
$\Sigma(c) = (S_{\text{scale}}(c), I_{\text{scale}}(c),
E_{\text{scale}}(c), \norm{\Delta}_{\max}(c))$
together with a healthy operating point
$(E,I,S)_{\text{healthy}}(c) \in [0,1]^3$. Each component carries
provenance metadata $(d, n, p, \rho)$ where $d$ is the ISO date of
measurement, $n$ is the corpus size, $p \in \{75, 90, 95, 99\}$ is
the percentile basis, and $\rho \in \{\text{measured}, \text{placeholder}, \text{derived}\}$
is the provenance tag.
\end{definition}

\subsection{Calibration Protocol}

For each class $c \in \mathcal{K}$, scale constants are measured by
the following protocol:

\begin{enumerate}
  \item \textbf{Corpus assembly.} Select all governance check-ins
    over a fixed window (the deployment uses 7 days) from agents
    $a$ with $k(a) = c$, restricted to sessions in which no
    \texttt{pause} or \texttt{reject} verdict was issued. This is
    the \emph{healthy} corpus for class $c$, denoted $\mathcal{H}_c$.
  \item \textbf{Raw quantity computation.} For each
    turn $t \in \mathcal{H}_c$, compute the unnormalized quantities
    $S_{\text{raw}}(t), I_{\text{raw}}(t), E_{\text{raw}}(t)$ using
    the highest-tier estimator available for that turn (logprob,
    multisample, or heuristic). Tag each measurement with its
    source.
  \item \textbf{Scale constants.} Set
    $S_{\text{scale}}(c) = \text{quantile}_{0.90}(\{S_{\text{raw}}(t)\}_{t \in \mathcal{H}_c})$,
    and similarly for $I_{\text{scale}}(c)$ and $E_{\text{scale}}(c)$.
    The 90th-percentile choice ensures typical healthy turns sit
    well inside $[0,1]$ after normalization.
  \item \textbf{Healthy operating point.} Set
    $(E,I,S)_{\text{healthy}}(c)$ to the median of the
    normalized $(E, I, S)$ values across $\mathcal{H}_c$.
  \item \textbf{Manifold scale.} Set
    $\norm{\Delta}_{\max}(c) = \text{quantile}_{0.95}(\norm{\Delta(a)}_2)$
    over $\mathcal{H}_c$. The 95th-percentile choice keeps the
    coherence floor near zero only for the most class-atypical
    states.
  \item \textbf{Provenance.} Stamp each constant with $d$
    (current ISO date), $n = |\mathcal{H}_c|$, $p$ (percentile
    used), and $\rho = \text{measured}$. Persist alongside the
    constant value.
  \item \textbf{Re-measurement cadence.} Re-run quarterly. Drift
    in any constant beyond a configurable threshold is itself a
    signal: the class's healthy envelope is shifting.
\end{enumerate}

\begin{remark}[Insufficient class population]
\label{rem:insufficient-class}
If $|\mathcal{H}_c|$ falls below a configured threshold
$N_{\min}$, the class-conditional measurement is skipped and the
class falls back to $\Sigma(c_{\text{default}})$. Production
deployments with very small class populations (e.g.,
\texttt{embodied} with a single Lumen instance) carry this
fallback explicitly until the population is large enough to
support an independent measurement. The fallback is itself
recorded with provenance $\rho = \text{placeholder}$ to
distinguish it from genuinely measured constants.
\end{remark}

\subsection{Fallback and Composition}

Two fallback regimes apply.

\paragraph{Single-instance classes.} For classes with only one
deployed agent (e.g., \texttt{embodied} for the present
deployment), the calibration protocol degenerates: the agent's own
history defines its healthy envelope, with no cross-instance
comparison possible. We treat this as a corner case where the
class is effectively the agent. The provenance stays
$\rho = \text{placeholder}$ until enough turns accumulate to
construct a meaningful percentile, after which it transitions to
$\rho = \text{measured}$. This is the operating mode for
\texttt{Lumen} during the framework's first six months of
deployment.

\paragraph{Composite tag classes.} An agent may carry multiple
tags (e.g., $\{\texttt{autonomous}, \texttt{persistent}\}$). The
class indicator $k(a)$ is a function of the agent's tag set, not
a single tag. We adopt the convention that more specific labels
(e.g., \texttt{Sentinel}) supersede tag-only classification when
both are available; absent a label, the tag set is treated as a
multi-key index into $\Sigma$. Classes that share a tag but
differ in label may share scale constants if the calibration data
shows their healthy envelopes are statistically indistinguishable;
subsequent work will report which tag-grouped classes met that
condition in the production corpus.

\begin{remark}[Cold-start: a bootstrapping window]
\label{rem:cold-start}
During the first $N_{\min}$ turns of a new class, the
fallback-to-$\Sigma(c_{\text{default}})$ behavior is, formally,
fleet-wide normalization applied to that single class. This is a
bounded-duration bootstrapping window, not a steady-state
commitment: the provenance tag $\rho = \text{placeholder}$ makes
the regime explicit in every check-in payload, and once the class
accumulates $N_{\min}$ healthy turns the measurement pipeline
transitions the provenance to \texttt{measured}. Operators
deploying a new class with known-different tempo or modality can
hand-author an initial $\Sigma(c)$ with provenance
\texttt{derived} to skip the window entirely.
\end{remark}

\subsection{What This Buys}

Two governance failure modes are eliminated by class-conditional
calibration that survive under fleet-wide normalization:

\begin{description}
\item[False incoherence.] Under fleet-wide
  $(E,I,S)_{\text{healthy}}$, an agent class with a genuinely
  high-entropy operating mode (creative/divergent) registers as
  permanently low-coherence. The verdict layer treats this as a
  governance signal. Class-conditional baselines correct this:
  the divergent agent's coherence is measured against the
  divergent class's own median, not against an analytical-class
  median.
\item[False resource alarms.] The resource-form energy
  $E = (\text{tokens}/s) / (\text{tokens}/s)_{\max}$ is meaningful
  only when the tokenizer and inference pipeline are comparable
  across the population. A Haiku-class agent operating at $80$
  tokens/s under a fleet-wide
  $(\text{tokens}/s)_{\max} = 200$ registers
  $E = 0.4$; an Opus-class agent at $40$ tokens/s registers
  $E = 0.2$, even though both are operating well within their
  respective envelopes. Class-conditional
  $E_{\text{scale}}$ corrects this.
\end{description}

The contribution of this section is the protocol, not the
specific numbers. The numbers belong to the empirical
appendix that lands with Phase 2 calibration; the protocol is a
design commitment that survives any specific fleet composition.

\begin{remark}[Calibration as evidence of precision miscalibration]
\label{rem:calibration-precision}
Under the FEP reading of \Cref{sec:grounding} ($E = -F$, with $F$
a variational free energy built from precision-weighted prediction
errors), precision is the inverse variance of the likelihood on
observations. Calibration deviation---the empirical mismatch
between an agent's stated confidence and the frequency with which
its claims are corroborated---is not literally identical to
precision deviation. It is \emph{evidence} of precision
miscalibration under one assumption: the agent's reported
confidence proxies the inverse-variance weighting its generative
model places on its predictions. Under that assumption, a
systematically overconfident agent is one that has placed
too-high precision on low-quality evidence, and the calibration
protocol above becomes an empirical handle on precision
miscalibration. When the assumption fails---e.g., reported
confidence is shaped by a constitutional statement rather than a
likelihood weight, or by a tokenized probability that is not the
agent's own precision---the protocol still measures something
operationally useful, but its FEP-native reading weakens. The
protocol stands on its own; the precision reading is a bridge to
the active-inference literature, not a claim of equivalence.
\end{remark}


% ═══════════════════════════════════════════════════════════
\section{Stochastic Extensions}
\label{sec:stochastic}

In production, the dynamics are subject to stochastic perturbations
from unpredictable task inputs. We model this as an It\^{o} SDE~\cite{ito1944}:
\begin{equation}
  d\eisv = f(\eisv, t)\,dt + G(\eisv, t)\,dW_t
  \label{eq:sde}
\end{equation}
where $W_t$ is a standard Wiener process and
$G(\eisv, t) \in \R^{4 \times 4}$ is the diffusion matrix.

\begin{theorem}[Mean-Square Contraction]
\label{thm:stochastic}
If the deterministic system is contracting with rate $\alpha_c$ under
metric $M$, and the diffusion satisfies
\begin{equation}
  \operatorname{tr}(G^\top M G) \leq \sigma_{\max}^2
  \label{eq:diffusion-bound}
\end{equation}
then the stochastic system satisfies:
\begin{multline}
  \mathbb{E}\bigl[\norm{\eisv_1(t) - \eisv_2(t)}_M^2\bigr] \leq \\
  e^{-2\alpha_c t}\norm{\eisv_1(0) - \eisv_2(0)}_M^2
  + \frac{\sigma_{\max}^2}{2\alpha_c}
  \label{eq:mean-square}
\end{multline}
The steady-state variance is bounded by
$\sigma_{\max}^2 / (2\alpha_c)$.
\end{theorem}

\begin{remark}[Verification for Production Noise Model]
In production, the diffusion matrix is taken as
$G = \operatorname{diag}(0, 0, \sigma, 0)$ with $\sigma = 0.02$
(additive noise on the entropy channel only). This is a modeling
choice chosen to be conservatively large relative to the entropy
fluctuations observable in the audit log under stationary task load;
an empirical trace-level estimate of $\sigma$ from per-agent $S_t$
time series is future work. The qualitative conclusion below (noise
an order of magnitude below typical signal excursions) tolerates
$\sigma$ varying by a factor of $2$ in either direction; a larger
departure would warrant re-derivation. For any diagonal metric
$M = \operatorname{diag}(m_E, m_I, m_S, m_V)$:
\[
  \operatorname{tr}(G^\top M G) = m_S \sigma^2
\]
Setting $m_S = 1$ (natural scaling), $\sigma_{\max}^2 = (0.02)^2 =
4 \times 10^{-4}$. With the certified contraction rate
$\alpha_c = 0.019$ (\Cref{app:contraction-proof}), the steady-state
variance bound from \eqref{eq:mean-square} is:
\[
  \frac{\sigma_{\max}^2}{2\alpha_c} = \frac{4 \times 10^{-4}}{0.038}
  \approx 0.011
\]
Using the tighter eigenvalue rate $\alpha_c \approx 0.13$
(Remark~\ref{rem:production-params}) would yield $1.5 \times 10^{-3}$,
but the conservative Gershgorin bound already shows the noise is
negligible. This is one order of magnitude smaller than typical EISV signal
excursions ($\sim 0.1$--$0.3$), confirming that stochastic noise does
not compromise governance decisions. The result follows directly from
Pham et~al.~\cite{pham2009stochastic}.
\end{remark}

The mean-square bound of \Cref{thm:stochastic} is a worst-case
contraction statement; a verdict-level false-alarm rate
$P_{\text{FA}} = \int_{\{\Phi(\eisv) < 0\}} p_\ast(\eisv)\, d\eisv$
requires the stationary distribution of the SDE and is deferred
to subsequent work. Related Fokker--Planck frameworks for LLM
drift~\cite{carson2025stochastic} give closed-form first-passage
times for one-dimensional severity processes.


% ═══════════════════════════════════════════════════════════
\section{Multi-Agent Network}
\label{sec:network}

When multiple agents operate simultaneously, their governance
states can be coupled through shared information. The deployed
system does not implement inter-agent state coupling---each agent's
EISV dynamics evolve independently---and the analysis below
therefore describes the theoretical behavior of a coupled
extension, not the current production mechanism. The synchronization
result of \Cref{thm:network} is used elsewhere in the paper only as
a within-class qualification on homogenization claims.

\subsection{Network Topology}

Consider $N$ agents with states $\eisv_i$, $i = 1, \ldots, N$,
connected by a graph $\mathcal{G} = (\mathcal{V}, \mathcal{E})$ with
Laplacian $L$. The coupled dynamics are:
\begin{equation}
  \dot{\eisv}_i = f(\eisv_i, t)
  + \epsilon_c \sum_{j \in \mathcal{N}_i} H(\eisv_j - \eisv_i)
  \label{eq:network}
\end{equation}
where $H \in \R^{4 \times 4}$ is a coupling matrix and
$\epsilon_c > 0$ is the coupling strength.

\begin{theorem}[Network Synchronization, Within-Class]
\label{thm:network}
Consider a coupled subnetwork in which every agent shares the same
isolated dynamics $f$ contracting with rate $\alpha_c$. The
subnetwork synchronizes (all such agents converge to a common
trajectory) provided
\begin{equation}
  \alpha_c > \lambda_{\max}(L) \cdot \lambda_{\max}(H)
  \label{eq:sync-condition}
\end{equation}
where $\lambda_{\max}$ denotes the largest eigenvalue.
\end{theorem}

\begin{remark}[Synchronization is a within-class property]
\Cref{thm:network} is a same-dynamics result and therefore applies
within an agent class as defined in \Cref{sec:class-structure}, not
across the heterogeneous fleet as a whole. Two agents with
class-conditional scale constants and healthy operating points that
differ between classes do not have ``the same'' isolated dynamics in
the sense the theorem requires; expecting them to converge to a
common trajectory is the homogenization failure mode discussed in
\Cref{sec:heterogeneity}. The relevant operational claim across
classes is not synchronization but coordinated divergence within each
class's own envelope.
\end{remark}


% ═══════════════════════════════════════════════════════════
\section{Experiments}
\label{sec:experiments}

This section reports a single production snapshot of the deployed
UNITARES system measured on February 20, 2026. All counts, tables, and
figures in this section are intended to describe that snapshot unless
explicitly stated otherwise. The purpose of the section is not to claim
full empirical validation of the grounded, class-conditional framework;
rather, it documents the operating characteristics of the deployed
system, the behavior of the existing dynamics in production, and the
practical motivation for the migration described in
\Cref{sec:migration}.

The results below should therefore be read as partial deployment
evidence. They support three narrower claims: that the framework is
deployed in a nontrivial production setting, that the observed fleet
dynamics are consistent with the qualitative behavior analyzed in
earlier sections, and that the migration and calibration problems are
real system-level concerns rather than hypothetical ones.

\subsection{Production Snapshot}

The UNITARES framework is deployed as a Model Context Protocol (MCP)
server that provides governance services to autonomous AI agents. Each
agent connects via standard MCP transport and receives governance
verdicts on every action. The system has been in continuous production
since December 2025.

\begin{table}[htbp]
\centering
\small
\caption{Production Deployment Summary}
\label{tab:deployment}
\begin{tabular}{@{}lr@{}}
\toprule
\textbf{Metric} & \textbf{Value} \\
\midrule
Registered agents & 903 \\
Active ($>5$ updates) & 75 \\
Max updates (single agent) & 48{,}668 \\
Audit events & 198{,}333 \\
Knowledge discoveries & 500 \\
Dialectic sessions & 66 \\
Deployment duration & 69 days \\
\bottomrule
\end{tabular}
\end{table}

The 903 registered agents include both long-running production agents
(5 agents with $>1{,}000$ updates, peak at 48{,}668) and shorter-lived
agents from development and testing cycles. The top two agents have
accumulated 48{,}668 and 44{,}529 updates respectively, providing
substantial trajectory data for convergence analysis.


\subsection{Observed Fleet Dynamics}

For the 75 active agents with $>5$ updates in the snapshot, the
observed equilibrium statistics are reported in
\Cref{tab:eisv-observed}. These are pooled across the agent
classes outlined in \Cref{sec:heterogeneity}; per-class values are
deferred to subsequent work, as discussed in
\Cref{sec:limitations-and-scope} and \Cref{sec:class-calibration}.

\begin{table}[htbp]
\centering
\small
\caption{Observed fleet-aggregate EISV equilibrium (75 active agents). Per-class breakdown deferred to subsequent work.}
\label{tab:eisv-observed}
\begin{tabular}{@{}lrrr@{}}
\toprule
\textbf{Variable} & \textbf{Mean} & \textbf{Std} & \textbf{Range} \\
\midrule
$E$ & 0.766 & 0.087 & $[0.60, 1.00]$ \\
$I$ & 0.877 & 0.111 & $[0.59, 1.00]$ \\
$S$ & 0.082 & 0.045 & $[0.01, 0.20]$ \\
$V$ & $-0.033$ & 0.028 & $[-0.089, 0.007]$ \\
$C(V)$ & 0.483 & 0.012 & $[0.46, 0.50]$ \\
\bottomrule
\end{tabular}
\end{table}

\begin{remark}[Reading fleet-aggregates honestly]
The pooled statistics in \Cref{tab:eisv-observed} mix an embodied
agent's sensor-driven state with session-bounded text-generation
agents and resident cron-driven janitors. The means and standard
deviations are real measurements and useful for cross-validation
against the dynamical predictions of \Cref{sec:dynamics}, but
they do not characterize any single agent class accurately. A
high-entropy divergent agent and a low-entropy convergent agent
contribute equally to the $S$ column; whether either is
``healthy'' is a class-conditional question that
\Cref{sec:class-calibration} formalizes and subsequent work
answers empirically.
\end{remark}

\Cref{tab:eisv-observed} summarizes the observed equilibrium.
These values are consistent with several qualitative predictions
of the dynamics:
\begin{enumerate}
\item \textbf{$I > E$:} Mean integrity (0.877) exceeds mean energy
  (0.766), confirming the system's conservative bias
  (\Cref{fig:ei-scatter}). The negative mean $V = -0.033$ (integrity
  surplus) is consistent with this.
\item \textbf{Low entropy:} Mean $S = 0.082$ with strong decay ($\mu
  = 0.8$) indicates effective uncertainty management.
\item \textbf{Coherence at 0.49:} The observed mean $C = 0.483$
  (\Cref{fig:coherence-hist}) confirms
  Remark~\ref{rem:operating-range}---the system honestly reports
  near-midpoint coherence rather than inflating the signal.
\item \textbf{$V$ confined to $[-0.1, 0.1]$:} 100\% of active agents
  have $V$ within $[-0.089, 0.007]$ (\Cref{fig:sv-scatter}), validating
  the damping analysis.
\end{enumerate}

\begin{figure}[t]
\centering
\includegraphics[width=\columnwidth]{fig1_ei_scatter}
\caption{Energy--Integrity scatter plot for 75 active agents, colored by
dynamical regime. Nearly all agents lie above the $I = E$ diagonal,
confirming the system's conservative bias ($\bar{V} < 0$). Point size
reflects log update count.}
\label{fig:ei-scatter}
\end{figure}

\begin{figure}[t]
\centering
\includegraphics[width=\columnwidth]{fig2_sv_scatter}
\caption{Entropy--Void operating region. All agents have $V \in [-0.089,
0.007]$, well within the $[-0.1, 0.1]$ predicted by damping analysis.
Entropy remains low ($S < 0.20$) due to strong decay $\mu = 0.8$.}
\label{fig:sv-scatter}
\end{figure}

\begin{figure}[t]
\centering
\includegraphics[width=\columnwidth]{fig3_coherence_hist}
\caption{Distribution of coherence $C(V)$ across active agents. The mean
$\bar{C} = 0.483$ lies near the theoretical midpoint ($C = 0.5$),
consistent with small $|V|$ and the honest-reporting principle.}
\label{fig:coherence-hist}
\end{figure}

\subsection{Dynamical Regimes}

Agents distribute across three dynamical regimes based on EISV
trajectory trends:

\begin{table}[htbp]
\centering
\small
\caption{Regime Distribution and EISV Profiles (Active Agents)}
\label{tab:regimes}
\begin{tabular}{@{}lrrrrr@{}}
\toprule
\textbf{Regime} & $n$ & $\bar{E}$ & $\bar{I}$ & $\bar{S}$ & $\bar{C}$ \\
\midrule
Convergence & 24 & .838 & .961 & .052 & .474 \\
Exploration & 22 & .766 & .926 & .090 & .477 \\
Divergence & 29 & .706 & .770 & .100 & .496 \\
\bottomrule
\end{tabular}
\end{table}

The regime profiles (\Cref{fig:regime-profiles}) show clear
separation: convergent agents have the highest integrity and lowest
entropy; divergent agents have the lowest integrity and highest
entropy. Coherence is slightly
higher in the divergent regime ($\bar{C} = 0.496$) because these
agents have smaller $|V|$, placing them nearer the coherence
function's midpoint.

\begin{figure}[t]
\centering
\includegraphics[width=\columnwidth]{fig5_regime_profiles}
\caption{Mean EISV component values by dynamical regime. Convergent
agents show the highest $I$ and lowest $S$; divergent agents show the
reverse pattern. The $|V|$ is small across all regimes, confirming
effective damping.}
\label{fig:regime-profiles}
\end{figure}

\subsection{Operational Artifacts: Knowledge Graph and Dialectic}
\label{sec:operational-artifacts}

The governance system has accumulated 500 knowledge graph discoveries
classified by severity: 238 low, 154 medium, 49 high, and 4 critical.
By type, the largest categories are insights (183), notes (125),
improvements (94), and bugs found (48).

The dialectic protocol---a structured thesis/antithesis/synthesis
process for resolving governance disputes---has conducted 66 sessions:
39 recovery sessions, 15 reviews, and 12 explorations. Of these,
21 reached resolution, while 45 terminated without agreement,
suggesting that the dialectic process is more effective for recovery
(where the system has clear signals) than for open-ended review.
The present deployment has no autonomous reviewer agents---the
resident fleet monitors, janitors, pattern-matchers, sync
daemons, and sensory processes, none of which carries a
thesis-response pathway---so every antithesis in the dataset
required human routing. The 45/66 non-resolution rate should
therefore be read as an upper bound under facilitator load rather
than as a protocol-level failure rate; a fleet with autonomous
reviewers would be expected to resolve more readily.

\paragraph{Retrieval infrastructure (April 2026).}
The knowledge graph's counts above describe \emph{storage}; whether
agents could \emph{locate} prior discoveries was a separate question.
A dogfooding review in mid-April 2026 found that semantic retrieval
over the discovery corpus was returning similarity scores near the
noise floor across diverse queries, tag filters were post-filter
rather than ranking signals, and an unlabeled low-threshold fallback
was returning random-looking results as if they were matches. The
retrieval layer was rebuilt in a five-phase sequence
(2026-04-20): an honest-failure-mode fix, a reproducible eval
harness, a modern embedder (\texttt{BAAI/bge-m3},
1024-dimensional), and hybrid retrieval via reciprocal rank fusion
over BM25 and dense scores. Cross-encoder reranking and 1-hop
typed-edge expansion were landed as opt-in modes; on the current
corpus size they showed no aggregate gain. The rebuild is
prerequisite to any empirical claim that relies on agents
\emph{using} prior discoveries rather than merely \emph{storing}
them; documentation and baselines live in the governance repository
under \texttt{docs/plans/2026-04-20-kg-retrieval-rebuild.md}.

\begin{figure}[t]
\centering
\includegraphics[width=\columnwidth]{fig6_discovery_types}
\caption{Knowledge graph discovery classification by type. Insights
and notes dominate, with bugs and improvements constituting the
actionable governance outputs.}
\label{fig:discovery-types}
\end{figure}

\subsection{Phase~2 Measurement: Per-Class Envelopes}
\label{sec:phase-2-measurement}

The Phase~2 calibration protocol of \Cref{sec:class-calibration}
was executed against a $30$-day healthy-regime slice of
\texttt{core.agent\_state} on April~18,~2026. \Cref{tab:per-class}
reports the resulting per-class healthy operating points and
manifold radii. The script that produced these values
(\texttt{scripts/calibrate\_class\_conditional.py}) ships with the
implementation; the values themselves are now committed in
\texttt{config/governance\_config.py} as the production
class-conditional constants, replacing the Phase~1 placeholders
for the manifold-coherence path.

\begin{table}[t]
\centering\small
\caption{Phase~2 measured per-class scale constants from a $30$-day
healthy-regime slice (April~18,~2026). $(E, I, S)_{\text{healthy}}$
is the median; $\norm{\Delta}_{\max}$ is the $95$th percentile of
state-space distance from each class's own healthy operating point.
Five classes with $N \geq 30$ qualified for measurement; classes
below the threshold fall back to the fleet-wide default.}
\label{tab:per-class}
\begin{tabular}{@{}lrrrrr@{}}
\toprule
\textbf{Class} & \textbf{N} & $E_h$ & $I_h$ & $S_h$ & $\norm{\Delta}_{\max}$ \\
\midrule
Lumen     & 7320 & 0.745 & 0.800 & 0.168 & 0.119 \\
default   & 2033 & 0.726 & 0.793 & 0.236 & 0.202 \\
Sentinel  & 1870 & 0.751 & 0.798 & 0.193 & 0.170 \\
Vigil     &  384 & 0.737 & 0.790 & 0.240 & 0.171 \\
Watcher   &  283 & 0.748 & 0.769 & 0.248 & 0.395 \\
\bottomrule
\end{tabular}
\end{table}

Two observations from \Cref{tab:per-class} are worth surfacing.

First, the per-class manifold radii span a $3.3\times$ range
(Lumen $0.119$ to Watcher $0.395$). The Phase~1 placeholder of
$\norm{\Delta}_{\max} = 1.8$ was therefore $4$ to $15\times$
larger than the empirically measured envelope of any class. A
fleet-wide normalization at the placeholder value would have
produced a manifold coherence saturating near $1.0$ for every
agent regardless of state---essentially a constant signal. This
is the quantitative motivation behind the homogenization argument
of \Cref{sec:heterogeneity}: a fleet-wide constant chosen to
cover the worst case (or chosen arbitrarily) loses discriminating
power within any class.

Second, the healthy operating points cluster more tightly on
$E$ and $I$ than on $S$. Across the five measured classes,
$E_h \in [0.726, 0.751]$ and $I_h \in [0.769, 0.800]$, but
$S_h \in [0.168, 0.248]$---a wider spread on entropy. This
matches the qualitative expectation that what makes classes
behaviorally distinct is largely how much uncertainty they
operate under: Lumen at $S_h = 0.17$ runs the most convergent;
Watcher at $S_h = 0.25$ runs more divergent because its work
(pattern-match over recently-edited code) is intrinsically
exploratory.

\Cref{tab:per-class} replaces the Phase~1 placeholder values for
the manifold-coherence path; the remaining
$S_{\text{scale}}, I_{\text{scale}}, E_{\text{scale}}$ constants
(used only by the not-yet-shipped tier-1 logprob path and
tier-2 multi-sample path) remain as Phase~1 placeholders.

\subsection{Verdict Counterfactual}
\label{sec:verdict-counterfactual}

The grounded manifold-coherence form (\Cref{eq:C-manifold}) takes
different inputs than the legacy operational form
(\Cref{eq:coherence}): the grounded form computes a
class-conditional distance from $(E,I,S)_{\text{healthy}}(c)$,
while the legacy form maps $V$ through $\tanh$. The choice of
coherence formula therefore has consequences at the gating layer,
not just on the raw value reported. We quantify this directly by
re-running the basin-classification step on a counterfactual in
which the grounded form replaces the legacy form on the same
production state vectors.

The basin classifier \texttt{classify\_basin} maps
$(E,I,S,V,C,R)$ to one of three regions---\texttt{high},
\texttt{boundary}, \texttt{low}---consumed by the production
dashboard and escalation logic alongside the
$\Phi$-thresholded verdict of \Cref{eq:verdict}. A
\emph{basin flip} is a row whose classifier output differs
between the legacy and grounded coherence inputs; it is the
observable downstream of which this counterfactual measures.

\textbf{Setup.} For each row in a $30$-day window of
\texttt{core.agent\_state} ($N = 13{,}310$ rows with computable
$E$), we compute $C_{\text{legacy}}$ (the value already stored in
the row, produced by the production server using
\Cref{eq:coherence}) and $C_{\text{grounded}}(c)$ (the manifold
form \Cref{eq:C-manifold} with class-conditional
$(E,I,S)_{\text{healthy}}(c)$ and $\norm{\Delta}_{\max}(c)$ from
\Cref{tab:per-class}, where $c$ is determined by the identity-tag
classifier of \Cref{sec:class-structure}). Each row is classified
into a basin twice using the same
\texttt{classify\_basin($E,I,S,V,C,\text{risk}$)} function used
in production. A flip is a row whose basin assignment differs.

\textbf{Result.} 3{,}844 of 13{,}310 rows ($28.9\%$) flip basin
assignment. \Cref{tab:counterfactual} reports per-class breakdown
and \Cref{tab:flip-directions} reports flip directions.

\begin{table}[t]
\centering\small
\caption{Verdict counterfactual: per-class basin flip rates when the
grounded class-conditional manifold coherence replaces the legacy
fleet-wide $\tanh$ form on the same state vectors.}
\label{tab:counterfactual}
\begin{tabular}{@{}lrrrr@{}}
\toprule
\textbf{Class} & \textbf{N} & \textbf{Flips} & \textbf{\%} & $|\Delta C|_{\max}$ \\
\midrule
Lumen     & 7889 & 2548 & 32.3 & 0.51 \\
default   & 2358 &  722 & 30.6 & 0.50 \\
Sentinel  & 2227 &  351 & 15.8 & 0.51 \\
Vigil     &  472 &  156 & 33.1 & 0.50 \\
Watcher   &  364 &   67 & 18.4 & 0.48 \\
\midrule
\textbf{Total} & \textbf{13310} & \textbf{3844} & \textbf{28.9} & \\
\bottomrule
\end{tabular}
\end{table}

\begin{table}[t]
\centering\small
\caption{Flip-direction breakdown by class. Predominance of
\texttt{high$\to$low} and \texttt{boundary$\to$low} indicates the
two formulations disagree asymmetrically: state vectors the
fleet-wide form classifies as healthy are more often reclassified
as breaching under the class-conditional form than vice-versa.}
\label{tab:flip-directions}
\begin{tabular}{@{}lrrr@{}}
\toprule
\textbf{Class} & \texttt{high$\to$low} & \texttt{bnd$\to$low} & \texttt{high$\to$bnd} \\
\midrule
Lumen     & 1169 & 758 & 621 \\
default   &  140 & 518 &  64 \\
Sentinel  &   35 & 289 &  27 \\
Vigil     &   70 &  70 &  16 \\
Watcher   &    0 &  67 &   0 \\
\bottomrule
\end{tabular}
\end{table}

\textbf{Interpretation.} The result is empirical evidence that the
choice of coherence formula has substantial governance-decision
consequences---the formula change is not a cosmetic relabeling
even when both quantities live in $[0,1]$. The flip rates of
$15$--$33\%$ are far above what could be attributed to formula
noise; they reflect the structural difference between a fleet-wide
$\tanh$ of $V$ and a class-conditional manifold distance from each
class's empirically measured healthy operating point.

The flip direction is informative. The dominant transition is into
the \texttt{low} basin: agents that the production server
classified as healthy (or borderline) under the fleet-wide form
are classified as breaching at least one bound under the
class-conditional grounded form. This is the gating-layer
footprint of the homogenization-failure-mode argument of
\Cref{sec:heterogeneity}: a fleet-wide constant spanning the
worst-case envelope is, by construction, permissive relative to
per-class envelopes, and the flip data quantifies how often that
permissiveness changes the verdict. The magnitude of the
consequence (29\% fleet-wide, up to 33\% per class) is empirical;
the directional asymmetry is a mechanical consequence of
replacing a fleet-wide envelope with per-class envelopes, not an
independent corroboration of the theory that motivated the
replacement.

We do \emph{not} claim that the grounded form is normatively
``correct'' in every flip---a flip into \texttt{low} would surface
a \texttt{guide} or \texttt{pause} verdict that the legacy form
suppressed, and whether that is desirable depends on the
agent. The contribution of this counterfactual is to establish
that the formula choice matters at the verdict level, not just at
the reported-value level. Calibration of which form should drive
governance decisions, and per-class tuning of the manifold radii
to match operator preferences, is left to subsequent work.

\textbf{Reproducibility.} The script that produces this table
(\texttt{scripts/verdict\_counterfactual.py}) is committed
alongside the calibration script. Running it against the
production database reproduces the per-class breakdown; running
with \texttt{-{}-csv} produces a row-level dump for further
analysis.

\subsection{Limits of the Snapshot}
\label{sec:limits-snapshot}

The production data has six limitations that bound the strength of
our empirical claims:
\begin{enumerate}
\item \textbf{Parameter evolution:} Dynamics parameters, governor
  thresholds, and the I-channel mode were refined iteratively during
  deployment. The reported EISV statistics are a snapshot across agents
  that experienced different parameter regimes. A controlled A/B
  comparison between logistic and linear modes has not been conducted.
\item \textbf{Governor recency:} The adaptive governor (CIRS~v2) was
  deployed recently; the audit log does not yet contain sufficient
  pre/post data for a rigorous before/after comparison.
\item \textbf{Population heterogeneity:} The 903 agents include
  long-running production agents, short-lived test agents, and agents
  from early development with since-changed parameters. We report
  results for the 75 agents with $>5$ updates to mitigate this, but
  acknowledge that even this filtered set spans multiple configuration
  epochs.
\item \textbf{Temporal seam:} Fleet-aggregate statistics
  (\Cref{tab:eisv-observed}) are from a February~20,~2026 snapshot;
  per-class envelopes (\Cref{sec:phase-2-measurement}) and the
  verdict counterfactual (\Cref{sec:verdict-counterfactual}) are
  from a 30-day window ending April~18,~2026, during which the
  population grew from 903 to 3{,}286 registered identities.
  The February~20 snapshot was taken against the v5 SQLite governance
  store, which was retired during the January~2026 PostgreSQL
  migration; the snapshot is preserved as historical record and is
  not reproducible against the current schema. Per-class envelopes
  and the verdict counterfactual were run against the post-migration
  PostgreSQL store and remain reproducible from current data.
\item \textbf{Identity-system maturity:} The measurement window
  includes material revisions to the identity-resolution layer that
  landed in mid-to-late April~2026---removal of name-claim lookup,
  adoption of UUID-direct dispatch, and receipt-based onboarding
  authentication. Class assignment is tag-based on the identity record
  (\Cref{sec:class-structure}), so an identity whose tags were
  inherited from a cached binding or from an archived predecessor
  may appear in a class whose envelope does not match its actual
  operational regime. We expect this to affect a bounded fraction of
  rows rather than the aggregate direction of the per-class envelopes,
  but we have not filtered the counterfactual to post-fix identities;
  a re-measurement on identity-clean data, once sufficient post-fix
  history has accumulated, is deferred to subsequent work.
\item \textbf{Retrieval infrastructure and write-path hygiene:}
  Two infrastructure defects predating the measurement window affect
  how the knowledge-graph counts (\Cref{tab:deployment}) and the
  dialectic resolution statistics should be read. First, until the
  April-2026 retrieval rebuild (\Cref{sec:operational-artifacts}),
  semantic retrieval scored at the noise floor and cross-agent
  learning was effectively gated by agents knowing a-priori where
  to look; fleet-learning claims that depend on agents
  \emph{using} prior discoveries are therefore operationally bounded
  by that pre-rebuild era. Second, the note-write handler
  silently auto-tagged every non-permanent discovery with
  \texttt{ephemeral} until PR \#54 (2026-04-19), which scheduled
  those notes for 7-day auto-archive; the tag-lifecycle counts in
  the archived tier thus contain both agent-chosen and
  handler-injected values in the pre-fix window. Neither defect
  affects the grounded EISV counterfactual of
  \Cref{sec:verdict-counterfactual} or the per-class envelopes of
  \Cref{sec:phase-2-measurement}, whose inputs are
  \texttt{core.agent\_state} rows rather than KG content; we flag
  them here because they bound operational claims
  about the shared-memory surface.
\end{enumerate}


% ═══════════════════════════════════════════════════════════
\section{Re-Grounding a Deployed Framework}
\label{sec:migration}

The counterfactual of \Cref{sec:verdict-counterfactual} was
computable because the framework was instrumented during a live
migration: grounded and legacy coherence were computed on every
production check-in over the measurement window, and governance
gating continued to read the legacy values so no verdict changed
during measurement. This section documents the mechanism that
made the comparison possible. We believe it generalizes to any
deployed agent framework whose coordinate semantics need to shift
while service continues.

\subsection{Three-Phase Dual-Compute}

The migration runs in three discrete phases:

\begin{description}
\item[Phase 1: Dual-compute.]
  Every check-in computes both the grounded value and the legacy
  value. Grounded values land in canonical
  $\{E, I, S, C\}$ slots; legacy values are preserved in
  $\{E_{\text{legacy}}, I_{\text{legacy}}, S_{\text{legacy}},
  C_{\text{legacy}}\}$. Source annotations on each grounded
  coordinate (\texttt{e\_source}, \texttt{i\_source},
  \texttt{s\_source}, \texttt{coherence\_source}) identify which
  computation tier produced each value (logprob / multisample /
  heuristic / resource / manifold / KL / FEP).
\item[Phase 2: Calibration.]
  Scale constants and class-conditional healthy operating points
  are measured against a reference corpus of healthy production
  agent-turns, partitioned by the identity tags of
  \Cref{sec:class-structure}. Divergence between grounded and
  legacy values is characterized; the counterfactual of
  \Cref{sec:verdict-counterfactual} is the first such
  characterization.
\item[Phase 3: Migration.]
  Governance gating (verdicts, basin classification, risk
  thresholds) switches from legacy to grounded values. The
  legacy fields remain readable for one minor version as a
  comparison surface, then are dropped.
\end{description}

The contribution of this section is not the three-phase shape
itself---which is standard for any deprecated-API migration---but
the mechanism by which Phase 1 preserves governance behavior while
exposing grounded values to downstream consumers.

\subsection{Pipeline-Ordering as a Behavior-Preservation Mechanism}
\label{sec:pipeline-ordering}

A naive reading of Phase 1 is that grounded values must equal
legacy values for behavior to be preserved. This is false in
general: the grounded coherence $C$ under \Cref{eq:C-manifold}
differs numerically from the legacy
$\tanh$-based coherence even when computed on identical EISV
inputs, because the formulas are different. Requiring numerical
equality would defeat the point of grounding.

The mechanism we use instead is pipeline-ordering. The server's
\texttt{process\_agent\_update} request pipeline runs in two
stages---which we will call the \emph{decision stage} (internally
named \texttt{phases} in the server code) and the \emph{payload
stage} (internally \texttt{enrichments}). The decision stage
computes governance state and issues the verdict; the payload
stage augments the response with secondary signals. Payload-stage
handlers are registered with integer priorities in the range
$[0, 100]$ (lower runs first). We register the grounding
computation as a payload-stage handler at priority $75$,
downstream of every handler that contributes to the verdict and
basin classification but upstream of presentation handlers such as
the human-readable summary and the broadcaster event.

The result: by the time the grounding handler runs, the verdict
has already been computed---using the legacy values produced by
the decision stage. The grounding handler then swaps grounded
values into the canonical
$\{E, I, S, C\}$ slots, copying the originals into
$\{E_{\text{legacy}}, \dots\}$. The response payload carries
grounded values, but the governance decision was made on legacy.
No gating code is touched; no consumer reading
$\{E, I, S, C\}$ sees a missing field; no decision changes during
Phase 1.

\begin{remark}[Scope of Phase 1]
$V$ is not dual-computed; its grounded reinterpretation as
accumulated free-energy debt (\Cref{eq:V-grounded}) is a
relabeling, not a new computation. Tier-1 (logprob) and tier-2
(multi-sample) grounding paths are interface stubs that raise
\texttt{NotImplementedError}, with the dispatcher falling through
to tier-3; inference-layer instrumentation for tier-1 lands
separately.
\end{remark}

\subsection{Generalization}

The pipeline-ordering mechanism is reusable for any deployed
agent framework facing a semantic-coordinate change while service
continues. The required structure is: (i) a request pipeline with
explicit ordering between decision-producing stages and
presentation-producing stages; (ii) a canonical payload slot for
the coordinate being re-grounded; (iii) identity (or session)
keying that allows audit consumers to tag rows by the migration
phase active when the row was produced.
Frameworks that conflate decision and presentation into a single
non-ordered handler must first separate the two concerns---a
separation worth doing independently of any migration, because it
makes governance decisions auditable and presentation choices
revisable in isolation.

As a concrete second example, consider an RLHF-governed deployment
whose gating signal is a reward-model score $r(x, y)$ and whose
operators wish to migrate to a calibrated reward $\tilde{r}(x, y)$
derived from a revised preference model without changing refusal or
escalation behavior during the transition. The pipeline structure
maps directly: the \emph{phases} stage computes the legacy $r$ and
issues the gate; the \emph{enrichments} stage computes $\tilde{r}$,
writes it to a canonical \texttt{reward} field, and copies $r$ to
\texttt{reward\_legacy}. Downstream dashboards and calibration
pipelines see $\tilde{r}$; refusal logic continues to read $r$.
The mechanism is identical; only the coordinate being re-grounded
changes.


% ═══════════════════════════════════════════════════════════
\section{Discussion}
\label{sec:discussion}

\subsection{Design Principle: Adjust Expectations, Not Measurements}
\label{sec:honest}

A core design principle of UNITARES is that observations are never
transformed to appear healthier. When production coherence averages
$C \approx 0.49$ rather than the ``expected'' $C \approx 0.7$, the
system reports 0.49. The phase-aware governor (\Cref{sec:governor})
addresses the gap by adjusting \emph{thresholds} (expectations), not
\emph{signals} (measurements).

This principle---\textbf{adjust expectations, not measurements}---has
three practical consequences:
\begin{enumerate}
\item \textbf{Calibration integrity:} Drift detection
  (\Cref{sec:ethical-drift}) compares raw signals against baselines.
  Inflating measurements would mask the very drift the system is
  designed to detect.
\item \textbf{Operator trust:} Dashboard values correspond to actual
  system state. An operator seeing $C = 0.49$ knows the agent is
  operating conservatively, not that a scaling factor was applied.
\item \textbf{Contraction validity:} The stability proofs in
  \Cref{sec:contraction} depend on the true state. Transforming
  measurements would invalidate the contraction analysis by
  introducing an unmodeled observation function.
\end{enumerate}

\subsection{Parameter Sensitivity}

The system's behavior is most sensitive to four parameters:

\paragraph{Convergence rate $\alpha$.}
Controls how quickly $E$ tracks $I$. Production uses
$\alpha = 0.42$. Larger $\alpha$ increases diagonal dominance in
the Jacobian (improving contraction) but makes the energy channel
less responsive to behavioral--epistemic drift.

\paragraph{Coherence coupling $\beta_I$.}
Determines how strongly coherence restores integrity. The theoretical
optimum for contraction ($\beta_I = 0.05$, minimizing off-diagonal
coupling) conflicts with the operational need for coherence feedback
($\beta_I = 0.30$, enabling meaningful recovery from entropy-driven
degradation). Production uses the larger value; contraction is
maintained because $\beta_I C'_{\max} = 0.15$ remains below the
diagonal terms.

\paragraph{Self-regulation $\gamma_I$.}
Determines the I-channel's intrinsic damping. In logistic mode,
$\gamma_I = 0.25$ yields $m_{\text{sat}} = -1.19$ at production
operating points (\Cref{app:saturation})---all agents are saturated,
relying on clamping rather than dynamics to bound $I$. In linear mode,
$\gamma_I = 0.169$ was tuned to give $I^* \approx 0.80$, well within
$(0,1)$, eliminating boundary dependence entirely.

\paragraph{Void damping $\delta$.}
At $\delta = 0.4$, the void channel has a time constant of
$\tau_V = 1/\delta = 2.5$ update cycles. This confines $V$ to
$[-0.1, 0.1]$ in production (\Cref{rem:operating-range}). Reducing
$\delta$ would widen the coherence operating range but slow the
system's response to $E$-$I$ imbalances.

\subsection{Verdict Sensitivity to the Objective Weights}
\label{sec:phi-sensitivity}

The scalar governance signal of \Cref{eq:phi} uses default weights
$w_E = w_I = w_S = w_V = w_\eta = 0.5$. The equal-weight choice is
a deployment default, not a derived optimum; an honest treatment
asks how sensitive the verdict surface is to weight perturbations.

At the verdict boundary $\Phi = 0.15$, the partial derivatives
against each weight are the corresponding state-dependent
contributions:
\begin{align*}
  \frac{\partial \Phi}{\partial w_E} &= E,
  & \frac{\partial \Phi}{\partial w_I} &= -(1-I),
  & \frac{\partial \Phi}{\partial w_S} &= -S, \\
  \frac{\partial \Phi}{\partial w_V} &= -\abs{V},
  & \frac{\partial \Phi}{\partial w_\eta} &= -\norm{\deta}^2.
\end{align*}
Evaluating at the production fleet-aggregate operating point
$(E, I, S, V) = (0.77, 0.88, 0.08, -0.03)$ with healthy
$\norm{\deta}^2 \lesssim 0.01$:
\begin{center}
\small
\setlength{\tabcolsep}{4pt}
\begin{tabular}{@{}lrl@{}}
\toprule
\textbf{Weight} & $\partial\Phi/\partial w_i$ & \textbf{Sign} \\
\midrule
$w_E$ & $+0.77$ & promotes proceed \\
$w_I$ & $-0.12$ & promotes proceed (small; $I$ high) \\
$w_S$ & $-0.08$ & promotes proceed (small; $S$ low) \\
$w_V$ & $-0.03$ & negligible \\
$w_\eta$ & $\lesssim 0.01$ & negligible \\
\bottomrule
\end{tabular}
\end{center}
The $w_E$ derivative dominates by an order of magnitude at healthy
operating points. This is structural rather than incidental: $E$
sits near $0.77$ in production while the penalty terms are all
near zero. The objective is therefore most sensitive to the
energy weight and least sensitive to the void/drift weights
\emph{in the healthy regime}. The ordering flips at degraded
states where $S$ rises and $(1-I)$ grows.

\paragraph{Threshold robustness band.}
Perturbing all five weights uniformly by $\pm \epsilon$ shifts
$\Phi$ by at most
$\epsilon \cdot (E + (1-I) + S + \abs{V} + \norm{\deta}^2)
\approx 1.01\epsilon$ in magnitude at the production operating
point. The verdict boundary $\Phi = 0.15$ is therefore robust to
uniform weight perturbations up to roughly $\pm 15\%$ before a
single-agent verdict at the boundary would flip. Non-uniform
perturbations are bounded by the largest single partial
derivative; a $+50\%$ perturbation to $w_E$ alone would shift
$\Phi$ by $+0.19$ at the operating point, which is well above
$0.15$ and would change boundary-case verdicts.

\paragraph{What this analysis does not do.}
The derivatives above characterize local sensitivity at a single
operating point; they do not establish how many verdicts would
flip under a specified weight perturbation, nor whether equal
weights are in any sense optimal. That result requires an
empirical counterfactual: rerun the audit log with perturbed
weights and count verdict flips, analogous to the grounded-form
counterfactual reported in \Cref{sec:experiments}. We defer the
weight-perturbation verdict-flip study to subsequent work with
the same infrastructure.

\subsection{Limitations}

\begin{enumerate}
\item The coherence function $\cv$ creates a nonlinear feedback
  loop whose gain depends on the operating point (maximum at
  $V = 0$, vanishing at $|V| \gg 0$). The contraction analysis
  uses $C'_{\max}$ as a worst-case bound, which is conservative.
\item The behavioral--epistemic drift vector's four components are currently
  equally weighted; adaptive weighting based on domain context
  is future work.
\item The multi-agent network analysis assumes fixed topology;
  production agent populations are dynamic with agents joining
  and leaving.
\item The PID governor gains ($K_p$, $K_i$, $K_d$) were tuned
  empirically; optimal gain selection via, e.g., Ziegler-Nichols
  or LQR methods is future work.
\item The empirical program of \Cref{sec:experiments} was computed
  over a window that overlaps with material identity-system revisions
  (\Cref{sec:limits-snapshot}, item~5). A replication of the per-class
  calibration (\Cref{sec:phase-2-measurement}) and the verdict
  counterfactual (\Cref{sec:verdict-counterfactual}) on identities
  that post-date those revisions is the natural validation path and
  is deferred to subsequent work.
\end{enumerate}

% ═══════════════════════════════════════════════════════════
\section{Conclusion}
\label{sec:conclusion}

This paper presented two structural reframings of UNITARES.
First, the EISV state vector is grounded information-theoretically:
each coordinate is a Shannon entropy, mutual information,
variational free energy, or KL/manifold distance, measured by a
tiered computation recipe and normalized by a provenance-tagged
scale constant. Second, fleet-wide normalization is rejected as
the wrong target; scale constants and healthy operating points
become functions of an agent class, keyed on existing identity
tags. Together these turn UNITARES from a phenomenological
dynamical system named after thermodynamic quantities into a
class-conditional information-theoretic governance framework
whose coordinates can be derived rather than guessed.

The pipeline-ordering migration mechanism
(\Cref{sec:pipeline-ordering}) is a methodological contribution
beyond UNITARES itself: any deployed agent framework that needs
to shift its semantic coordinates while preserving live governance
behavior can apply the same separation between decision-producing
and presentation-producing pipeline stages. We argue this
separation is worth doing on its own merits, independent of any
specific migration.

Stability of the dynamics is preserved under the grounded
reformulation, with full proof in \Cref{app:contraction-proof}.
Production deployment statistics from 903~agents over 69~days
appear in \Cref{sec:experiments}; \Cref{sec:phase-2-measurement}
reports a first class-conditional measurement for the
manifold-coherence path, and
\Cref{sec:verdict-counterfactual} shows that $28.9\%$ of basin
assignments flip under the grounded counterfactual. The remaining
scale constants and higher-tier estimators are deferred to
subsequent work following the measurement protocol of
\Cref{sec:class-calibration}.

The framework's adaptive governor (\Cref{sec:governor}),
behavioral--epistemic drift vector (\Cref{sec:ethical-drift}),
and multi-agent extensions (\Cref{sec:network}) form the rest of
the governance surface alongside these two commitments.

The current verdict surface---\texttt{proceed}, \texttt{guide},
\texttt{pause}, \texttt{reject}---is threshold-triggered from the
CIRS v2 control law (\Cref{sec:governor}), not derived from the
minimization of an expected free energy. A principled extension
under active inference~\cite{parr2022active,dacosta2020active}
would recast governance interventions as policy selections that
minimize expected free energy over future trajectories; the current
framework monitors free-energy proxies and intervenes via
phase-conditioned thresholds. We do not claim the existing verdicts
are active-inference actions, and the reframing is future work.

A cleaner identification of the deferred step: Tier-1 computation
of $-F$ (\Cref{sec:grounding-quantities}) requires an explicit
generative model $p(z, o)$ of the agent. For closed-weights frontier
agents, it is not obvious what that generative model \emph{is} in
the strict variational sense. The weights are not directly
accessible; a tokenized posterior over next tokens is a tractable
but coarse surrogate; a semantic generative model over tool-use
trajectories is richer but speculative. Which choice, if any,
supports a faithful $-F$ estimator is an open methodological
problem, not an engineering deferral. The natural empirical test
is whether a candidate $\hat{F}$ adds predictive value over the
current EISV proxies; that horse race is left to subsequent work.

As autonomous AI agents become more prevalent, governance
frameworks must accommodate the heterogeneity of the populations
they govern---embodied creatures, persistent observers,
session-bounded assistants, ephemeral parsers---without
collapsing them to a homogeneous distribution. UNITARES v6
presents one structural commitment to that requirement.


% ═══════════════════════════════════════════════════════════
% REFERENCES
% ═══════════════════════════════════════════════════════════
\bibliographystyle{plainnat}

\begin{thebibliography}{99}

\bibitem[Lohmiller and Slotine(1998)]{lohmiller1998contraction}
W.~Lohmiller and J.-J.~E. Slotine.
\newblock On contraction analysis for non-linear systems.
\newblock \emph{Automatica}, 34(6):683--696, 1998.

\bibitem[Friston(2010)]{friston2010free}
K.~Friston.
\newblock The free-energy principle: a unified brain theory?
\newblock \emph{Nature Reviews Neuroscience}, 11(2):127--138, 2010.

\bibitem[Da Costa et~al.(2020)]{dacosta2020active}
L.~Da Costa, T.~Parr, N.~Sajid, S.~Veselic, V.~Neacsu, and K.~Friston.
\newblock Active inference on discrete state-spaces: A synthesis.
\newblock \emph{Journal of Mathematical Psychology}, 99:102447, 2020.

\bibitem[Sajid et~al.(2021)]{sajid2021active}
N.~Sajid, P.~J. Ball, T.~Parr, and K.~J. Friston.
\newblock Active inference: demystified and compared.
\newblock \emph{Neural Computation}, 33(3):674--712, 2021.

\bibitem[Smith et~al.(2022)]{smith2022step}
R.~Smith, K.~J. Friston, and C.~J. Whyte.
\newblock A step-by-step tutorial on active inference and its application to empirical data.
\newblock \emph{Journal of Mathematical Psychology}, 107:102632, 2022.

\bibitem[Parr et~al.(2022)]{parr2022active}
T.~Parr, G.~Pezzulo, and K.~J. Friston.
\newblock \emph{Active Inference: The Free Energy Principle in Mind, Brain, and Behavior}.
\newblock MIT Press, 2022.

\bibitem[Shannon(1948)]{shannon1948}
C.~E. Shannon.
\newblock A mathematical theory of communication.
\newblock \emph{Bell System Technical Journal}, 27(3):379--423, 1948.

\bibitem[Slotine(2003)]{slotine2003modular}
J.-J.~E. Slotine.
\newblock Modular stability tools for distributed computation and control.
\newblock \emph{International Journal of Adaptive Control and Signal Processing},
  17(6):397--416, 2003.

\bibitem[Guo et~al.(2017)]{guo2017calibration}
C.~Guo, G.~Pleiss, Y.~Sun, and K.~Q. Weinberger.
\newblock On calibration of modern neural networks.
\newblock In \emph{ICML}, pages 1321--1330, 2017.

\bibitem[Niculescu-Mizil and Caruana(2005)]{niculescu2005predicting}
A.~Niculescu-Mizil and R.~Caruana.
\newblock Predicting good probabilities with supervised learning.
\newblock In \emph{ICML}, pages 625--632, 2005.

\bibitem[\AA{}str\"{o}m and Murray(2008)]{astrom2008feedback}
K.~J. \AA{}str\"{o}m and R.~M. Murray.
\newblock \emph{Feedback Systems: An Introduction for Scientists and Engineers}.
\newblock Princeton University Press, 2008.

\bibitem[Bai et~al.(2022)]{bai2022constitutional}
Y.~Bai, S.~Kadavath, S.~Kundu, et~al.
\newblock Constitutional {AI}: Harmlessness from {AI} feedback.
\newblock \emph{arXiv preprint arXiv:2212.08073}, 2022.

\bibitem[Ouyang et~al.(2022)]{ouyang2022training}
L.~Ouyang, J.~Wu, X.~Jiang, et~al.
\newblock Training language models to follow instructions with human feedback.
\newblock In \emph{NeurIPS}, 2022.

\bibitem[Rebedea et~al.(2023)]{rebedea2023nemo}
T.~Rebedea, R.~Dinu, M.~Sreedhar, C.~Parisien, and J.~Cohen.
\newblock {NeMo} {G}uardrails: A toolkit for controllable and safe {LLM}
  applications with programmable rails.
\newblock In \emph{EMNLP (Demo)}, 2023.

\bibitem[Ravindran(2025)]{ravindran2025moral}
S.~Ravindran.
\newblock A predictive framework for {AI} value alignment and drift prevention.
\newblock \emph{arXiv preprint arXiv:2510.04073}, 2025.

\bibitem[Kuhn et~al.(2023)]{kuhn2023semantic}
L.~Kuhn, Y.~Gal, and S.~Farquhar.
\newblock Semantic uncertainty: Linguistic invariances for uncertainty
  estimation in natural language generation.
\newblock In \emph{ICLR}, 2023.
\newblock arXiv:2302.09664.

\bibitem[Wang et~al.(2025a)]{wang2025mi9}
C.~L. Wang, T.~Singhal, A.~Kelkar, and J.~Tuo.
\newblock {MI9} -- agent intelligence protocol: Runtime governance for agentic {AI}
  systems.
\newblock \emph{arXiv preprint arXiv:2508.03858}, 2025.

\bibitem[Rath(2026)]{rath2026agentdrift}
A.~Rath.
\newblock Agent drift: Quantifying behavioral degradation in multi-agent {LLM}
  systems over extended interactions.
\newblock \emph{arXiv preprint arXiv:2601.04170}, 2026.

\bibitem[Carson(2025)]{carson2025stochastic}
J.~D. Carson.
\newblock A stochastic dynamical theory of {LLM} self-adversariality: Modeling
  severity drift as a critical process.
\newblock \emph{arXiv preprint arXiv:2501.16783}, 2025.

\bibitem[Ames et~al.(2019)]{ames2019control}
A.~D. Ames, S.~Coogan, M.~Egerstedt, G.~Notomista, K.~Sreenath, and
  P.~Tabuada.
\newblock Control barrier functions: Theory and applications.
\newblock In \emph{European Control Conference (ECC)}, pages 3420--3431, 2019.

\bibitem[Pham et~al.(2009)]{pham2009stochastic}
Q.-C. Pham, N.~Tabareau, and J.-J.~E. Slotine.
\newblock A contraction theory approach to stochastic incremental stability.
\newblock \emph{IEEE Transactions on Automatic Control}, 54(4):816--820, 2009.

\bibitem[It\^{o}(1944)]{ito1944}
K.~It\^{o}.
\newblock Stochastic integral.
\newblock \emph{Proceedings of the Imperial Academy}, 20(8):519--524, 1944.

\bibitem[Wang et~al.(2025b)]{wang2025probguard}
H.~Wang, C.~M. Poskitt, J.~Wei, and J.~Sun.
\newblock {ProbGuard}: Probabilistic runtime monitoring for {LLM} agent
  safety.
\newblock \emph{arXiv preprint arXiv:2508.00500}, 2025.

\end{thebibliography}


% ═══════════════════════════════════════════════════════════
% APPENDIX
% ═══════════════════════════════════════════════════════════
\appendix

\section{Parameter Tables}
\label{app:parameters}

\begin{table}[htbp]
\centering
\footnotesize
\caption{EISV Dynamics Parameters (Production Values)}
\label{tab:params}
\begin{tabular}{@{}llrl@{}}
\toprule
\textbf{Parameter} & \textbf{Symbol} & \textbf{Val.} & \textbf{Role} \\
\midrule
\multicolumn{4}{@{}l}{\textit{Energy channel}} \\
Convergence & $\alpha$ & 0.42 & $E$ tracks $I$ \\
Entropy erosion & $\beta_E$ & 0.10 & $S$ erodes $E$ \\
Drift excitation & $\gamma_E$ & 0.05 & Drift raises $E$ \\
\midrule
\multicolumn{4}{@{}l}{\textit{Information channel}} \\
Entropy degrad. & $k$ & 0.10 & $S$ degrades $I$ \\
Coherence boost & $\beta_I$ & 0.30 & $C(V)$ restores $I$ \\
Self-regulation & $\gamma_I$ & 0.25 & Logistic damping \\
Linear mode & $\gamma_I^{\text{lin}}$ & 0.169 & Linear damping \\
\midrule
\multicolumn{4}{@{}l}{\textit{Entropy channel}} \\
Natural decay & $\mu$ & 0.80 & $S$ relaxation \\
Drift injection & $\lambda_{1}$ & 0.30 & Base (adaptive) \\
Coherence red. & $\lambda_{2}$ & 0.05 & $C(V)$ reduces $S$ \\
Complexity & $\beta_c$ & 0.15 & Task $\to$ entropy \\
Noise & $\sigma$ & 0.02 & Stochastic term \\
\midrule
\multicolumn{4}{@{}l}{\textit{Void channel}} \\
Coupling & $\kappa$ & 0.30 & $E$-$I$ gap \\
Damping & $\delta$ & 0.40 & $V$ relaxation \\
\midrule
\multicolumn{4}{@{}l}{\textit{Coherence}} \\
Maximum & $C_{\max}$ & 1.0 & Upper bound \\
Sensitivity & $C_1$ & 1.0 & Tanh slope \\
\bottomrule
\end{tabular}
\end{table}

\begin{table}[htbp]
\centering
\footnotesize
\caption{CIRS v2 Adaptive Governor Parameters}
\label{tab:governor}
\begin{tabular}{@{}llrl@{}}
\toprule
\textbf{Parameter} & \textbf{Symbol} & \textbf{Val.} & \textbf{Role} \\
\midrule
Proportional & $K_p$ & 0.05 & Gentle corr. \\
Integral & $K_i$ & 0.005 & Drift corr. \\
Derivative & $K_d$ & 0.10 & Damping \\
Integral limit & $I_{\max}$ & 0.10 & Wind-up cap \\
\midrule
\multicolumn{4}{@{}l}{\textit{Reference points}} \\
Expl.\ $\tau$ ref & $r_\tau^{(e)}$ & 0.35 & Forgiving \\
Expl.\ $\beta$ ref & $r_\beta^{(e)}$ & 0.55 & Forgiving \\
Intg.\ $\tau$ ref & $r_\tau^{(i)}$ & 0.40 & Strict \\
Intg.\ $\beta$ ref & $r_\beta^{(i)}$ & 0.60 & Strict \\
\midrule
\multicolumn{4}{@{}l}{\textit{Safety bounds}} \\
$\tau$ floor/ceil & $\tau$ & 0.25/0.75 & Hard bounds \\
$\beta$ floor/ceil & $\beta$ & 0.20/0.70 & Hard bounds \\
\bottomrule
\end{tabular}
\end{table}


\begin{table}[htbp]
\centering
\footnotesize
\caption{Behavioral--Epistemic Drift Vector Parameters}
\label{tab:drift}
\setlength{\tabcolsep}{4pt}
\begin{tabular}{@{}llrl@{}}
\toprule
\textbf{Parameter} & \textbf{Symbol} & \textbf{Val.} & \textbf{Role} \\
\midrule
\multicolumn{4}{@{}l}{\textit{Baseline tracking}} \\
EMA smoothing & $\alpha_{\text{ema}}$ & 0.1 & Update rate \\
Warmup steps & $n_w$ & 5 & Drift suppression \\
\midrule
\multicolumn{4}{@{}l}{\textit{Stability measurement}} \\
Decision win. & $w_s$ & 20 & Consistency \\
Transition pen. & --- & /flip & Instability \\
\midrule
\multicolumn{4}{@{}l}{\textit{Coupling to dynamics}} \\
Energy excit. & $\gamma_E$ & 0.05 & Drift $\to dE$ \\
Entropy inj. & $\lambda_1$ & .05--.20 & Drift $\to dS$ \\
\bottomrule
\end{tabular}
\end{table}


\begin{table}[htbp]
\centering
\footnotesize
\caption{Grounding scale constants (\Cref{sec:scale-constants}). The
manifold-coherence path is now measured class-conditionally in
\Cref{tab:per-class}; the remaining constants below remain
Phase~1 placeholders pending tier-1/tier-2 calibration work.}
\label{tab:scale-constants}
\setlength{\tabcolsep}{4pt}
\begin{tabular}{@{}lllll@{}}
\toprule
\textbf{Constant} & \textbf{Sym.} & \textbf{Val.} & \textbf{Pctl} & \textbf{Prov.} \\
\midrule
Entropy & $S_{\text{scale}}$ & 3.0 & --- & placeholder \\
MI & $I_{\text{scale}}$ & 2.0 & --- & placeholder \\
Free energy & $E_{\text{scale}}$ & 1.0 & --- & placeholder \\
Manifold rad. & $\norm{\Delta}_{\max}(c)$ & per $c$ & 95 & measured \\
Tok./sec env. & $(\text{t/s})_{\max}$ & 200 & --- & placeholder \\
\bottomrule
\end{tabular}
\end{table}

The Pctl column records the percentile basis once the constant
has been measured against a healthy corpus (per
\Cref{sec:class-calibration}); the Provenance column is one of
\{\texttt{placeholder}, \texttt{measured}, \texttt{derived}\}
per \Cref{def:scale-map}. The $\norm{\Delta}_{\max}(c)$ row is now
measured class-conditionally in \Cref{tab:per-class}; the remaining
entries are still \texttt{placeholder} scaffolding awaiting the
tier-1/tier-2 calibration work. A later revision can restore the
date and corpus-size columns once those remaining constants are
measured.


\section{Proof of Contraction (Full)}
\label{app:contraction-proof}

This appendix presents the full theorem statement and proof of
EISV stability summarized in \Cref{sec:contraction}, together with
the equilibrium characterization. Both were demoted from the v5
main body to this appendix; the v6 main body retains only the
result statement.

\subsection*{Contraction Framework}

\begin{definition}[Contraction~\cite{lohmiller1998contraction}]
A system $\dot{\eisv} = f(\eisv, t)$ is \emph{contracting} with
rate $\alpha > 0$ in a metric $M(\eisv, t) \succ 0$ if the
generalized Jacobian satisfies
\begin{equation}
  F = \left(\dot{M} + M\frac{\partial f}{\partial \eisv}
  + \left(\frac{\partial f}{\partial \eisv}\right)^\top M\right)
  \preceq -2\alpha M
  \label{eq:contraction-condition}
\end{equation}
for all $\eisv \in \mathcal{X}$, $t \geq 0$. This guarantees that
all trajectories converge exponentially to a unique trajectory at
rate $e^{-\alpha t}$.
\end{definition}

\subsection*{Jacobian of the EISV System}

Define $J = \partial f / \partial \eisv$. For the governing
equations \eqref{eq:dE}--\eqref{eq:dV}:
\begin{equation}
J = \begin{pmatrix}
-\alpha - \beta_E S & \alpha & -\beta_E E & 0 \\
0 & -\gamma_I & -k & \beta_I C'(V) \\
0 & 0 & -\mu & -\lambda_2 C'(V) \\
\kappa & -\kappa & 0 & -\delta
\end{pmatrix}
\label{eq:jacobian}
\end{equation}
where $C'(V) = \frac{\partial C}{\partial V} =
\frac{C_{\max} C_1}{2}\operatorname{sech}^2(C_1 V)$ and the
$\deta$-dependent terms are treated as exogenous forcing.

\subsection*{Theorem and Proof}

\begin{theorem}[EISV Contraction]
\label{thm:contraction}
The EISV system with linear information dynamics
($g_I(I) = \gamma_I I$) is contracting with rate
\begin{equation}
  \alpha_c = \min\left\{
    \alpha + \beta_E S_{\min},\;
    \gamma_I,\;
    \mu,\;
    \delta
  \right\} - \epsilon(M)
  \label{eq:contraction-rate}
\end{equation}
where $\epsilon(M) \geq 0$ accounts for off-diagonal coupling in
the chosen metric $M \succ 0$. With the diagonal metric
$M = \operatorname{diag}(m_E, m_I, m_S, m_V)$, a sufficient
condition for $\alpha_c > 0$ is
\begin{equation}
  \min\{\alpha, \gamma_I, \mu, \delta\} >
  \max\!\bigl\{
    \kappa\!\sqrt{\tfrac{m_V}{m_E}},\,
    \beta_I C'_{\max}\!\sqrt{\tfrac{m_I}{m_V}}
  \bigr\}
  \label{eq:contraction-sufficient}
\end{equation}
where $C'_{\max} = C_{\max}C_1/2$.
\end{theorem}

\begin{remark}[Production Parameters]
\label{rem:production-params}
With production values ($\alpha = 0.42$, $\gamma_I = 0.25$,
$\mu = 0.8$, $\delta = 0.4$, $\kappa = 0.3$, $\beta_I = 0.3$,
$C'_{\max} = 0.5$), the diagonal terms are
$\{0.42, 0.25, 0.8, 0.4\}$ and the coupling bound is dominated by
$\kappa = 0.3$ and $\beta_I C'_{\max} = 0.15$. The Gershgorin
analysis below with diagonal metric
$M = \operatorname{diag}(0.1, 0.2, 1.0, 0.08)$ yields a certified
contraction rate $\alpha_c = 0.019$. This is conservative:
computing the eigenvalues of $(MJ + J^\top\! M)/(2M)$ directly
gives $\alpha_c \approx 0.13$, and the spectral abscissa of the
Jacobian $J$ at equilibrium is $-0.15$. The Gershgorin bound is
tight on Row~2 (the $I$-channel) because the $m_E \alpha$
cross-term competes with the weak $\gamma_I = 0.25$ diagonal; the
actual system converges roughly $8\times$ faster than the
certified bound.
\end{remark}

\subsection*{Equilibrium Characterization}

Setting \eqref{eq:dE}--\eqref{eq:dV} to zero and solving:

\begin{proposition}[Equilibrium]
\label{prop:equilibrium}
Under mild conditions ($\alpha > 0$, $\mu > 0$, $\delta > 0$,
$\gamma_I > 0$), the system admits a unique equilibrium
$\eisv^* = (E^*, I^*, S^*, V^*)$ satisfying
\begin{align}
  V^* &= \frac{\kappa(E^* - I^*)}{\delta} \label{eq:Vstar} \\
  E^* &= \frac{\alpha I^*}{(\alpha + \beta_E S^*)} + O(\norm{\deta}^2)
    \label{eq:Estar} \\
  I^* &= \frac{\beta_I C(V^*)}{\gamma_I + kS^*/\beta_I C(V^*)}
    \label{eq:Istar}
\end{align}
The fixed point is globally exponentially stable by
\Cref{thm:contraction}.
\end{proposition}

\subsection*{Full Construction}

We prove \Cref{thm:contraction} by constructing a diagonal metric
$M = \operatorname{diag}(m_E, m_I, m_S, m_V) \succ 0$ and showing
that the generalized Jacobian condition \eqref{eq:contraction-condition}
holds throughout the state space $\mathcal{X}$.

\subsection*{Step 1: Symmetric Part of $MJ + J^\top M$}

For diagonal $M$, the matrix $F = MJ + J^\top M$ has entries:
\begin{align*}
  F_{ii} &= 2 m_i J_{ii} \\
  F_{ij} &= m_i J_{ij} + m_j J_{ji}, \quad i \neq j
\end{align*}

The diagonal entries of $F$ (using linear $g_I$ mode) are:
\begin{align}
  F_{11} &= -2m_E(\alpha + \beta_E S) \leq -2m_E \alpha
    \label{eq:F11} \\
  F_{22} &= -2m_I \gamma_I
    \label{eq:F22} \\
  F_{33} &= -2m_S \mu
    \label{eq:F33} \\
  F_{44} &= -2m_V \delta
    \label{eq:F44}
\end{align}
where \eqref{eq:F11} uses $S \geq S_{\min} \geq 0$.

The nonzero off-diagonal entries are:
\begin{align}
  F_{12} = F_{21} &= m_E \alpha \label{eq:F12} \\
  F_{13} &= -m_E \beta_E E \leq 0 \label{eq:F13} \\
  F_{23} &= -m_I k \label{eq:F23} \\
  F_{24} &= m_I \beta_I C'(V) \label{eq:F24} \\
  F_{34} &= -m_V \lambda_2 C'(V) \label{eq:F34} \\
  F_{14} = F_{41} &= m_V \kappa \label{eq:F14}
\end{align}
where $F_{13}$ and $F_{34}$ are negative (aiding contraction)
and are bounded conservatively by their absolute values in the
Gershgorin analysis below.

Note: For diagonal $M$, $F_{12} = F_{21} = m_E J_{12} + m_I J_{21}
= m_E \alpha$.

\subsection*{Step 2: Gershgorin Circle Theorem}

For $F + 2\alpha_c M \preceq 0$, it suffices by Gershgorin that
for each row~$i$:
\begin{equation}
  F_{ii} + 2\alpha_c m_i + \!\sum_{j \neq i}\! |F_{ij}| \leq 0
  \label{eq:gershgorin}
\end{equation}

Bounding $|F_{13}| \leq m_E\beta_E$ conservatively (its sign aids
contraction), the four row conditions are:

\paragraph{Row 1 ($E$):}
\begin{equation}
  m_E(-\alpha + \alpha_c + \beta_E) + m_V\kappa \leq 0
  \label{eq:gersh-row1}
\end{equation}

\paragraph{Row 2 ($I$):}
\begin{equation}
  \begin{split}
  m_I(-2\gamma_I + 2\alpha_c + k &+ \beta_I C'_{\max}) \\
  &+ m_E\alpha \leq 0
  \end{split}
  \label{eq:gersh-row2}
\end{equation}

\paragraph{Row 3 ($S$):}
\begin{equation}
  \begin{split}
  -2m_S\mu + 2\alpha_c m_S &+ m_E\beta_E \\
  &+ m_I k + m_V\lambda_2 C'_{\max} \leq 0
  \end{split}
  \label{eq:gersh-row3}
\end{equation}

\paragraph{Row 4 ($V$):}
\begin{equation}
  \begin{split}
  m_V(-2\delta + 2\alpha_c + \kappa &+ \lambda_2 C'_{\max}) \\
  &+ m_I\beta_I C'_{\max} \leq 0
  \end{split}
  \label{eq:gersh-row4}
\end{equation}

\subsection*{Step 3: Metric Construction}

We seek $m_E, m_I, m_S, m_V > 0$ and $\alpha_c > 0$ satisfying
\eqref{eq:gersh-row1}--\eqref{eq:gersh-row4}. A natural choice is to
set the metric weights inversely proportional to the coupling they
receive:
\begin{equation}
  m_E = 1, \quad m_I = \frac{\alpha}{2\gamma_I}, \quad
  m_S = 1, \quad m_V = \frac{\alpha}{2\delta}
\end{equation}

\subsection*{Step 4: Numerical Verification}

Production parameters: $\alpha = 0.42$, $\beta_E = 0.10$,
$\beta_I = 0.30$, $\gamma_I = 0.25$, $k = 0.10$, $\mu = 0.80$,
$\kappa = 0.30$, $\delta = 0.40$, $\lambda_2 = 0.05$,
$C'_{\max} = 0.5$.

\paragraph{Initial attempt.}
The natural metric $m_E = 1$, $m_I = 0.84$, $m_S = 1$,
$m_V = 0.525$ fails Row~2 ($+0.21 + 1.68\alpha_c > 0$).
Reducing $m_I = 0.2$, $m_E = 0.1$ makes Row~2 the binding
constraint:
\[
  -0.008 + 0.4\alpha_c \leq 0
  \;\implies\; \alpha_c < 0.02
\]
We set $\alpha_c = 0.019$ (strictly less) to ensure all
inequalities are strict.

\paragraph{Verification.}
With $M = \operatorname{diag}(0.1, 0.2, 1.0, 0.08)$ and
$\alpha_c = 0.019$:

\smallskip
\begin{center}
\footnotesize
\begin{tabular}{@{}lrrl@{}}
\toprule
\textbf{Row} & \textbf{LHS} & \textbf{$< 0$?} & \\
\midrule
1\,($E$) & $-0.006$ & yes & \checkmark \\
2\,($I$) & $-0.0004$ & yes & \checkmark \\
3\,($S$) & $-1.529$ & yes & \checkmark \\
4\,($V$) & $-0.005$ & yes & \checkmark \\
\bottomrule
\end{tabular}
\end{center}

\smallskip
Row~4 initially failed with $m_V = 0.05$ ($+0.008$); adjusting
to $m_V = 0.08$ resolves it ($-0.005$) while Row~1 remains
satisfied ($-0.006$). Row~2 now has strict margin ($-0.0004$).

\medskip
\textbf{Result:} With $M = \operatorname{diag}(0.1, 0.2, 1.0,
0.08)$, the Gershgorin analysis certifies $\alpha_c = 0.019$.
This is conservative---the eigenvalue-based bound is
$\alpha_c \approx 0.13$ and the spectral abscissa of~$J$ at
equilibrium is $-0.15$
(Remark~\ref{rem:production-params}). The key result is
$\alpha_c > 0$, guaranteeing global exponential
convergence. \qed


\section{I-Channel Self-Regulation: Linear vs.\ Logistic}
\label{app:saturation}

\Cref{sec:equations} specifies the linear self-regulation term
$g_I(I) = \gamma_I I$. This appendix records the logistic
alternative $g_I^{\log}(I) = \gamma_I I(1-I)$, its saturation
pathology under production forcing, and the empirical basis for
choosing the linear form.

\paragraph{Logistic form and saturation criterion.}
Under the logistic self-regulation, the $I$ subsystem
($dI/dt = 0$) has two equilibria
\begin{equation}
  I^*_{\pm} = \tfrac{1}{2} \pm \tfrac{1}{2}\sqrt{1 - 4A/\gamma_I},
  \qquad A = -kS + \beta_I C(V).
  \label{eq:I-equilibria}
\end{equation}

\begin{proposition}[Saturation criterion]
\label{prop:saturation}
The logistic $I$-dynamics has a boundary saturation pathology
when $A > \gamma_I/4$. In this regime $I \to 1$ monotonically and
the system loses its ability to distinguish integrity levels. The
saturation margin
\begin{equation}
  m_{\text{sat}} = 1 - 4A/\gamma_I
  \label{eq:saturation-margin}
\end{equation}
is positive in safe operation and non-positive under saturation.
\end{proposition}

\paragraph{Linear form.}
The linear self-regulation $g_I^{\text{lin}}(I) = \gamma_I I$
yields the unique equilibrium
\begin{equation}
  I^* = (-kS + \beta_I C(V)) / \gamma_I
  \label{eq:I-star-linear}
\end{equation}
which is always in $(0,1)$ when $\gamma_I$ exceeds the maximum
forcing. Under the linear form the Jacobian entry
$J_{22} = -\gamma_I$ is constant; under the logistic form
$J_{22} = -\gamma_I(1 - 2I)$, which changes sign at $I = 0.5$ and
approaches zero at the boundaries---precisely where contraction
is most needed. Production uses the linear form with
$\gamma_I = 0.169$, giving $I^* \approx 0.80$.

\paragraph{Empirical saturation under the logistic form.}
The saturation margin was computed for all $75$ active agents
under the logistic mode with parameters $\gamma_I = 0.25$,
$k = 0.10$, $\beta_I = 0.30$:

\begin{center}
\small
\begin{tabular}{@{}lrr@{}}
\toprule
\textbf{Regime} & \textbf{Mean $m_{\text{sat}}$} & \textbf{\% Sat.} \\
\midrule
Convergence ($n=24$) & $-1.45$ & 100\% \\
Exploration ($n=22$) & $-1.18$ & 100\% \\
Divergence ($n=29$) & $-1.10$ & 100\% \\
\midrule
All active ($n=75$) & $-1.23$ & 100\% \\
\bottomrule
\end{tabular}
\end{center}

Every active agent has a negative saturation margin under the
logistic form: the forcing term $A$ exceeds $\gamma_I/4 = 0.0625$
across the fleet. The saturation criterion
(\Cref{prop:saturation}) holds regardless of the agent's
configuration history, since it depends only on the relationship
$A > \gamma_I/4$.


\section{Boundary Clamping and Non-Smoothness}
\label{app:boundary-clamping}

The EISV state is clamped to $\mathcal{X}$ at each integration step,
introducing non-smoothness at domain boundaries. Three mechanisms
keep the dynamics well-behaved despite this:

\paragraph{Interior operation.}
The epistemic humility floor ($S_{\min} = 0.001$) and the contraction
dynamics keep trajectories in the interior of $\mathcal{X}$ during
normal operation. In production, only $19$ of $903$ agents (2.1\%)
contact the $I = 1$ boundary, all in the \texttt{CONVERGENCE} regime
with few updates; interior operation is the norm.

\paragraph{Projected dynamical systems.}
When boundary contact does occur, the clamping operation implements
a projection onto the constraint set $\mathcal{X}$. Projected
dynamical systems preserve contraction properties when the
constraint set is convex~\cite{lohmiller1998contraction}, which
$\mathcal{X}$ (a hyperrectangle) is.

\paragraph{Pre-emptive tightening.}
The saturation diagnostic of \Cref{app:saturation} monitors the
I-channel's distance from the boundary. When $m_{\text{sat}} < 0.1$
the system logs a warning and the adaptive governor tightens
thresholds pre-emptively, preventing boundary contact rather than
relying on clamping to maintain invariance.


\end{document}
